% Compile with pdfLaTeX (two passes); bibliography and all content are embedded.
\documentclass[11pt]{article}
\usepackage[margin=1in]{geometry}
\usepackage{amsmath,amssymb,amsthm,mathtools,booktabs,longtable,array}
\usepackage[T1]{fontenc}
\usepackage{lmodern,microtype}
\usepackage[colorlinks=true,allcolors=blue,bookmarksnumbered=true]{hyperref}
\hypersetup{pdftitle={Conditional Gaussian Preparation with Retained Archives},
 pdfauthor={Jeremy Rodgers},pdfsubject={Conditional preparation, retained archives and quantum instruments}}
\numberwithin{equation}{section}
\newtheorem{theorem}{Theorem}[section]
\newtheorem{lemma}[theorem]{Lemma}
\newtheorem{proposition}[theorem]{Proposition}
\newtheorem{corollary}[theorem]{Corollary}
\theoremstyle{definition}
\newtheorem{definition}[theorem]{Definition}
\newtheorem{assumption}[theorem]{Assumption}
\newtheorem{example}[theorem]{Example}
\theoremstyle{remark}
\newtheorem{remark}[theorem]{Remark}
\newcommand{\TV}{\operatorname{TV}}
\newcommand{\Dtr}{\operatorname{D}}
\newcommand{\tr}{\operatorname{tr}}
\newcommand{\id}{\operatorname{id}}
\newcommand{\dd}{\,\mathrm d}
\setlength{\emergencystretch}{2em}
\title{Conditional Gaussian Preparation\protect\\with Retained Archives}
\author{Jeremy Rodgers\thanks{Website: \href{https://everythingequation.com}{everythingequation.com}}\\Independent Researcher}
\date{9 October 2026\\[0.4em]
\normalsize\href{https://doi.org/10.5281/zenodo.23259560}{doi:10.5281/zenodo.23259560}}
\begin{document}
\maketitle
\begin{abstract}
We construct an exact smooth split--copy--return protocol for a Gaussian writer and a finite bank of retained position archives. Its complete guidance endpoint map is an area-preserving deformation of a baker map. At every finite separation, a singular archive population, or an exactly retained copy of an initial archive coordinate, can preserve the writer's full fine information: writer-only regularity does not imply conditional preparation. We prove a positive alternative using weak physical relative scores of the complete original conditional law. Sectional bounded-variation estimates and explicitly charged actual tails give quantitative averaged conditional readiness jointly with all final archive coordinates, without a density cap or finite quantile Fisher information. A finite $102$-coordinate example gives readiness below $1.0052\times10^{-5}$ and, using one common initial box, the same ceiling for a subsequent labelled projective instrument with an arbitrary finite internal reference. A distinct receiver copies the earlier writer event and retains its sign under the specified holding dynamics. Uniform full conditional writer-score bounds also yield finite-resource existence as the archive bank grows. The results use a declared controlled canonical-current parent; independent material control and preparation-law warrant remain separate hypotheses.
\end{abstract}
\clearpage
\tableofcontents
\newpage
\section{The preparation question}\label{sec:introduction}
Can reversible dynamics prepare a position variable for a quantum measurement while keeping the physical coordinates that recorded its earlier history? A marginal convergence statement does not answer this question. A downstream instrument may interact with an archive that still resolves the writer's initial position. Conversely, retaining an archive need not prevent conditional preparation when the complete initial law has appropriate regularity: information can remain in the archive while the writer's dyadic remainder approaches a uniform conditional law.

We construct and analyse this distinction for a finite prescribed harmonic model. A smooth split, copy and return protocol gives an exact two-coordinate map. Its Jacobian is one in reference cumulative-distribution coordinates, and it converges to a baker comparison as the Gaussian branch separation grows. The physical map, rather than a branchwise rigid-translation surrogate, is used throughout. Every archive coordinate and internal memory is retained. The source populations may be correlated and different from their wave densities.

The principal preparation theorem bounds the total variation between the final joint law and the product of a ready writer with its \emph{actual} nuisance marginal. Its assumptions concern weak physical relative scores of the complete original conditional density. Sectional bounded variation estimates remove the density cap and avoid assuming finite Fisher information after a Gaussian quantile transformation. An explicit finite example uses $100$ preparation archives and a separate measurement receiver. We then prove an input-uniform projective-instrument estimate for an arbitrary finite internal reference, including a distinct physical copy of the earlier writer event and its subsequent retention. The common initial-box formulation charges the original tail once.

Three features delimit the positive result. First, the law regularity is conditional on every genuinely retained external variable; a good writer marginal does not suffice. Second, the return factorizes the writer's \emph{wave}, whereas statistical readiness requires a separate estimate on the transported actual law. Third, the instrument is formulated in a declared canonical-current model with prescribed controls. The theorem does not supply the independent physical warrant for its initial statistical class or its control architecture.

\subsection{Relation to established and prior work}\label{sec:literature}
Equivariance and conditional wavefunctions are standard elements of the quantum-equilibrium analysis of D\"urr, Goldstein and Zangh\`i~\cite{DurrGoldsteinZanghi1992,DurrGoldsteinZanghi2004}. Equivariant uniqueness under locality assumptions, as studied by Goldstein and Struyve~\cite{GoldsteinStruyve2007}, addresses a different question from finite preparation of an admitted nonequilibrium class. Valentini and Westman's numerical relaxation studies~\cite{ValentiniWestman2005} concern coarse-grained evolution. Here the target metric keeps complete fine archive coordinates, and the initial density need not equal the quantum reference.

Triangular probability transformations~\cite{Rosenblatt1952}, bounded-variation methods for expanding maps~\cite{LasotaYorke1973}, Gaussian information inequalities~\cite{Gross1975}, and reversible computation with retained information~\cite{Bennett1973} are established antecedents. The conditional state at an actual configuration is distinguished from a reduced density matrix as in~\cite{DurrEtAl2005Density}. Our labelled state is compared with the standard projective instrument of quantum measurement theory~\cite{DaviesLewis1970,Ozawa1984}, using elementary pure-state trace-distance identities~\cite{FuchsVanDeGraaf1999}. This comparison is not called a diamond-norm bound: a linear completely positive actual map for arbitrary nonequilibrium laws has not been assumed or proved.

The author's earlier return-holonomy manuscript~\cite{RodgersReturn2026} develops an invariant-return kernel, pointwise mixing, finite-library obstructions and a retained-memory entropy identity with inverse echo. The control-consistency manuscript~\cite{RodgersControl2026} addresses uniqueness under a different statistical premise. These are contextual predecessors; neither uniqueness theorem is a premise of the deterministic preparation estimates proved here. Appendix~\ref{app:regular} gives a self-contained total-variation formulation of retained memory and a complementary regular-basin argument. The contribution of the present paper is the explicit nonsingular harmonic realization, the fine-coordinate obstruction, the full conditional score/moment preparation estimate, and their compatible one-use instrument assembly. The cited October research editions are distinct from the preserved September editions and from the new companion manuscripts discussed in Appendix~\ref{app:scope}.

\section{Laws, configurations and error metrics}\label{sec:notation}
All configuration coordinates in this paper belong to one specified effective parent. The writer is a real coordinate $q$ of width $\sigma$; preparation archives are $z_1,\ldots,z_n$ of widths $a_j$; a distinct, initially untouched receiver is $z_{n+1}$. A finite collection of two-state internal registers carries coherent branch information. Internal states have no independently sampled spin coordinate in this model. $X$ denotes any further genuinely retained external information. Conditional laws given $X=x$ are integrated against the \emph{actual} law of $X$.

Write $\phi(t)=(2\pi)^{-1/2}e^{-t^2/2}$ and $\Phi(t)=\int_{-\infty}^t\phi(s)\,ds$. Standardized coordinates are $\xi_0=q/\sigma$ and $\xi_j=z_j/a_j$. The product reference density $\Gamma_d(\xi)=\prod_{i=0}^{d-1}\phi(\xi_i)$ is fixed by the factored ready wave. The actual conditional density is $p_x=\Gamma_d f_x$ when it exists. This identity defines a relative density; it does not assign $f_x=1$. Analysis quantiles are $u_i=\Phi(\xi_i)$; no physical controller is assumed to evaluate a cumulative distribution function.

For equal-mass finite positive measures define $\TV(\mu,\nu)=\frac12\lVert\mu-\nu\rVert_1$. For probability measures this equals the supremum of event-probability differences. For density operators use $\Dtr(A,B)=\frac12\lVert A-B\rVert_1$, also for subnormalized operators of equal trace. Classical total variation contracts under a common Markov kernel, including a deterministic pushforward; trace distance contracts under a common quantum channel. The classical--quantum comparisons below also use the triangle inequality for mixtures of positive operators. If $N$ is the full nuisance configuration, the readiness functional is
\begin{equation}\label{eq:readiness-definition}
 \mathcal R(\mu)=\TV\bigl(\mu_{u,N},\lambda_{(0,1)}\otimes\mu_N\bigr).
\end{equation}
It measures averaged conditional readiness and keeps $\mu_N$ unchanged. Bounds on it imply bounds for every bounded downstream response using both variables, but not uniform conditional bounds after arbitrarily rare postselection. We apply the same definition to subprobabilities; both measures then have their common original mass.

The physical relative scores are weak derivatives $\partial_i f_x=f_x s_{i,x}$, with $J_i=\mathbb E_X\int s_{i,x}^2p_x\,d\xi$. Equivalently, the square root of the relative density has the corresponding weighted weak Sobolev derivative where the usual finite-information equivalence applies. Ordinary almost-everywhere derivatives that miss jumps are not sufficient. A twentieth-moment budget below means $\sum_i\mathbb E|\xi_i|^{20}$, not a radial moment or quantum energy. Directional bounded variation is always integrated over the remaining coordinates and $X$; when a density is cut to a physical box, both boundary faces are included.

Under the complete canonical current, $j_i=(\hbar/m_i)\operatorname{Im}(\Psi^*\partial_i\Psi)$, summed over every internal component, the reference density is $\rho=\lVert\Psi\rVert^2$ and the actual velocity is $j_i/\rho$. The active Gaussian reference density is strictly positive and each active velocity is a convex combination of finite trap velocities. These active trajectories exist uniquely at every finite entrance point over the finite protocol. In the concrete Gaussian bank the full positional density is also strictly positive. An external variable $X$ may instead be a fixed conditioning label, with no assigned wave amplitude. Any additional physical spectator must have its stipulated conservative holding flow, with the complete wave and conditional internal state defined on a set of full actual probability. All physical coordinates that move belong to the transported configuration. This specifies the flow domain used here; no pointwise uniqueness assertion for arbitrary singular quantum models is imported.

\newcommand{\pTV}{\operatorname{TV}}
\newcommand{\pVar}{\operatorname{Var}}
\newcommand{\pLaw}{\operatorname{Law}}
\newcommand{\pGam}{\Gamma}
\newcommand{\pUnif}{\mathsf U}

\section{A smooth copy--return map on the complete configuration}
\label{p1p:sec:map}

Write $\phi(x)=(2\pi)^{-1/2}e^{-x^2/2}$ and
$\Phi(x)=\int_{-\infty}^x\phi(y)\,dy$. A writer coordinate $q$ and a distinct
archive coordinate $z$ have prescribed oscillator ground-state densities
$g_\sigma(q)=\sigma^{-1}\phi(q/\sigma)$ and
$h_a(z)=a^{-1}\phi(z/a)$. Their quantum widths satisfy
$\sigma^2=\hbar/(2m\omega)$ and $a^2=\hbar/(2M\Omega)$.
The internal two-state degree of freedom is a fibre, with no additional
sampled configuration coordinate. The balanced entrance wave is
\begin{equation}\label{p1p:eq:stock}
 \Psi_{\rm in}(q,z)=\sqrt{g_\sigma(q)h_a(z)}
 (|0\rangle+|1\rangle)/\sqrt2.
\end{equation}
Equation~\eqref{p1p:eq:stock} specifies the quantum stock. The actual
configuration law will be a separate input to the preparation theorem.

The prescribed Hamiltonian admits exact translated packets. If $\varphi$
is the ground envelope of an oscillator and $d(t)$ is a translation, then
\begin{equation}\label{p1p:eq:translated}
 \psi_d(q,t)=\varphi(q-d(t))
 \exp\{im\dot d(t)q/\hbar+i\theta_d(t)\},\qquad
 \dot\theta_d=-\omega/2-m(\dot d)^2/(2\hbar),
\end{equation}
solves the Schr\"odinger equation with potential
\begin{equation}\label{p1p:eq:potential}
 V_d(q,t)=\tfrac12m\omega^2(q-d(t))^2-m\ddot d(t)q.
\end{equation}
This follows by differentiating~\eqref{p1p:eq:translated}: the transport
term $-i\hbar\dot d\varphi'$ cancels the cross kinetic term, the inertial
linear term cancels $-m\ddot d q$, and the remaining scalar terms give the
displayed phase. The density and current are exactly
$g_\sigma(q-d)$ and $\dot d\,g_\sigma(q-d)$.
For an infinitely differentiable control, define
\begin{equation}\label{eq:smooth-step}
 s(t)=\frac{\displaystyle\int_0^t e^{-1/[r(1-r)]}\,dr}
 {\displaystyle\int_0^1 e^{-1/[r(1-r)]}\,dr},\qquad 0<t<1,
\end{equation}
and extend it by $0$ for $t\leq0$ and $1$ for $t\geq1$.
Every positive-order derivative vanishes at both joins. On a stage of
duration $T>0$, take $d(t)=d_0+(d_1-d_0)s(t/T)$ after translating the
time origin. The resulting potential is smooth in time and space.
The polynomial $35t^4-84t^5+70t^6-20t^7$, joined to constants, is an
alternative $C^3$ schedule giving a $C^1$ time-dependent potential.
Both schedules give the same endpoint map below. Every finite translation
has finite prescribed coefficients; no uniform bound as $T\downarrow0$
is asserted.

The protocol consists of a writer stage, a copy stage, and a return
stage. The writer traps move conditionally on $|0\rangle,|1\rangle$ to
$-D\sigma,+D\sigma$. With the writer packets stationary, the archive
traps move to $-Aa,+Aa$. With the archive packets stationary, the writer
traps return to the origin. Orthogonality of the internal columns makes
the full positional density the sum of their positive densities.
The inactive coordinate has zero current in each stage: during copying
every writer envelope is real with only a $q$-independent phase; during
return every archive envelope is real with only a $z$-independent phase.
The active velocity is a convex combination of the two trap velocities.
It is smooth, locally Lipschitz and bounded on each finite time interval,
so every finite entrance point has a unique global stage trajectory.

Hereafter the intermediate physical positions are standardized, so $q$
means $q/\sigma$ and $z$ means $z/a$. Let
\begin{align}
 F_D(q)&=\tfrac12\{\Phi(q+D)+\Phi(q-D)\},&
 w_D(q)&=\frac{1}{1+e^{2Dq}},\label{p1p:eq:cdfs}\\
 H_{w,A}(z)&=w\Phi(z+A)+(1-w)\Phi(z-A),&
 w_A(z)&=\frac{1}{1+e^{2Az}},\nonumber\\
 M_A(z)&=\tfrac12\{\Phi(z+A)+\Phi(z-A)\}.&&\nonumber
\end{align}
In this section $M_A$ denotes a CDF; its derivative is the archive
reference density. The analysis coordinates are $u=\Phi(q_0)$ and
$v=\Phi(z_0)$. They are changes of mathematical variables; the control
does not evaluate a CDF.

\begin{theorem}[Complete Gaussian endpoint map]\label{p1p:thm:map}
For $D,A>0$, define
\begin{align}
 q&=F_D^{-1}(u),& z&=H_{w_D(q),A}^{-1}(v),\label{p1p:eq:map}\\
 r&=w_A(z)\Phi(q+D)+(1-w_A(z))\Phi(q-D),& b&=M_A(z).\nonumber
\end{align}
The exact guidance endpoint map in entrance and final reference
coordinates is $T_{D,A}(u,v)=(r,b)$. It is a smooth bijection of $(0,1)^2$
with determinant one. Its inverse is
\begin{align}
 z&=M_A^{-1}(b),&
 q&=[w_A(z)\Phi(\,\cdot+D)+(1-w_A(z))\Phi(\,\cdot-D)]^{-1}(r),
 \label{p1p:eq:inverse}\\
 u&=F_D(q),&v&=H_{w_D(q),A}(z).\nonumber
\end{align}
At the endpoint the writer quantum wave is again a common ground
envelope, factorized from the coherent archive--internal state.
\end{theorem}
\begin{proof}
In a one-dimensional continuity equation with vanishing current at
$-\infty$, the CDF $F_t$ obeys
$\partial_tF_t=-j_t$. Along $\dot q=j_t/\rho_t$,
$dF_t(q(t))/dt=0$. Applying this identity first to the writer gives
$q=F_D^{-1}(u)$. During copying $q$ is fixed, so the conditional archive
weights are the fixed numbers $w_D(q),1-w_D(q)$ and its conditional CDF
is conserved. During return $z$ is fixed; the conditional writer weights
are $w_A(z),1-w_A(z)$. Its final CDF is the ready CDF, giving $r$.
The final archive reference CDF is $M_A$, giving $b$.

Strict positivity gives every inverse in~\eqref{p1p:eq:inverse} and the
implicit-function theorem gives smoothness. Put
\[
 \rho(q,z)=\tfrac12\{\phi(q+D)\phi(z+A)+\phi(q-D)\phi(z-A)\}.
\]
The triangular input transformation $(q,z)\mapsto(u,v)$ has Jacobian
$F_D'(q)\partial_zH_{w_D(q),A}(z)=\rho(q,z)$.
The triangular output transformation $(q,z)\mapsto(r,b)$ has Jacobian
\[
 M_A'(z)\{w_A(z)\phi(q+D)+(1-w_A(z))\phi(q-D)\}=\rho(q,z).
\]
Their ratio is one. Finally, with the phases from each stage retained,
the wave is
\[
 \sqrt{g_\sigma(q)}\,
 \frac{e^{i\theta_0}\sqrt{h_a(z+Aa)}|0\rangle+
 e^{i\theta_1}\sqrt{h_a(z-Aa)}|1\rangle}{\sqrt2}.
\]
This proves factorization as well as the claimed complete map.
\end{proof}

To repeat the map, supply fresh internal qubits in $|+\rangle$ and fresh
archive oscillators, and retain the old qubits as internal memories.
An exact smooth internal swap also rotates the archive holding
projectors. If $V(t)$ is a spatially constant interpolation from the
identity to SWAP, set
\begin{equation}\label{p1p:eq:swap}
 H(t)=V(t)H_{\rm hold}V(t)^\dagger+i\hbar\dot V(t)V(t)^\dagger.
\end{equation}
Direct differentiation shows that the wave is $V(t)$ times the holding
wave. Thus all positional densities and currents are preserved during
the swap. One choice is
$V(t)=\exp[-is(t)\pi(I-\mathrm{SWAP})/2]$.
After each swap, all previously written archive currents are zero and
future operations fix their coordinates. The complete $j$th cycle is
therefore $T_{D,A}$ on the current writer and the $j$th unused archive,
with every other configuration coordinate fixed. This exact holding and
conjugated-control specification is part of the prescribed parent.

\subsection{Finite separation and retained fine information}

The comparison map is the invertible, area-preserving baker map
\begin{equation}\label{p1p:eq:baker}
 B(u,v)=(2u-s,(v+s)/2),\qquad s=\lfloor2u\rfloor,
\end{equation}
defined off its null branch cut. It is a comparison of complete maps,
not an instruction to translate actual points rigidly according to a
branch sign at finite overlap.

\begin{proposition}[A finite fine-archive obstruction]\label{p1p:prop:channel}
Let the initial writer quantile be uniform and the initial archive
quantile be a fixed $v_0\in(0,1)$. At every finite $D,A>0$,
\[
 \pTV\bigl(\pLaw(r,b),\pUnif\otimes\pLaw(b)\bigr)=1.
\]
The same conclusion holds conditionally if a regular initial archive
quantile is copied exactly into separately retained memory.
\end{proposition}
\begin{proof}
The copy equation is
$v_0=w_D(q)\Phi(z+A)+(1-w_D(q))\Phi(z-A)$.
Its implicit derivative is
\[
 \frac{dz}{dq}=-\frac{w_D'(q)[\Phi(z+A)-\Phi(z-A)]}
 {\partial_zH_{w_D(q),A}(z)}>0.
\]
Indeed $w_D'<0$ and both remaining factors are positive. From $(v_0,z)$
one recovers
\[
 w=\frac{v_0-\Phi(z-A)}{\Phi(z+A)-\Phi(z-A)},\qquad
 q=\frac{1}{2D}\log\frac{1-w}{w}.
\]
Thus $b$ determines $z,q,r$. The actual joint measure is carried by a
measurable graph, whereas its own $b$ marginal times a nonatomic uniform
$r$ law gives that graph measure zero. This gives total variation one.
Conditioning on an exactly retained initial $v$ gives the same proof.
\end{proof}

\begin{lemma}[Uniform sectional map comparison]\label{p1p:lem:separation}
Choose $c>0$, $D>c$, $0<\eta<1/2$, and $k>0$ such that
\[
 \gamma=e^{-2Dc}\leq\eta/2,\qquad \Phi(-k)\leq\eta/2,
 \qquad A>k.
\]
Put $a_D=\Phi(-(D+c))$, $t_A=\Phi(-(2A-k))$ and
$\rho_A=e^{-2A(A-k)}$, and assume $\gamma+t_A<\eta$.
Then $K=B^{-1}T_{D,A}$ has coordinatewise displacement at most
\begin{equation}\label{p1p:eq:delta-beta}
 \delta=\max\{(a_D+\rho_A)/2,\gamma+t_A\}
\end{equation}
outside a set of area at most
\begin{equation}\label{p1p:eq:beta}
 \beta=2\eta+\Phi(-(D-c)).
\end{equation}
On this good set the archive sign agrees with the writer sign just
before copying. The bounds hold on every fixed spectator fibre.
\end{lemma}
\begin{proof}
Use the good set $|F_D^{-1}(u)|\geq c$ and
$v\in[\eta,1-\eta]$. For $q\leq-c$, $1-w_D(q)\leq\gamma$.
The copy equation gives
\[
 \Phi(z+A)\leq\frac{1-\eta}{1-\gamma}\leq1-\eta/2,
 \qquad z\leq-A+k.
\]
Consequently $1-w_A(z)\leq\rho_A$ and $\Phi(z-A)\leq t_A$.
Since $2u=\Phi(q+D)+\Phi(q-D)$,
\[
 |r-2u|\leq a_D+\rho_A,
 \qquad 0\leq b-v/2\leq(\gamma+t_A)/2.
\]
The latter inequality and $v\leq1-\eta$ imply $b<1/2$;
also $z<0$. The left inverse-baker formula $(r,b)\mapsto(r/2,2b)$
therefore gives~\eqref{p1p:eq:delta-beta}. Reflection proves the right
case. The excluded writer-band length is exactly
$\Phi(-(D-c))-\Phi(-(D+c))$, bounded as in~\eqref{p1p:eq:beta}.
The two archive ends have total length $2\eta$.
\end{proof}

\section{Complete-law preparation with every archive retained}
\label{p1p:sec:preparation}

Let $X$ be arbitrary genuinely retained external information, on a
standard Borel space with its actual law $\mu_X$. During the present
protocol $X$ is fixed. Include in the entrance vector the writer
$U\in(0,1)$, all $n$ archive inputs $V_1,\ldots,V_n$, and any untouched
future positional sources $W$. Conditional on $X=x$, suppose their
complete joint law has a density $f_x$ relative to product Lebesgue
measure $\lambda$. The archive inputs may be arbitrarily correlated
with one another, the writer and $X$. All variations below are variations
of this full conditional density and are averaged using $\mu_X$.
In particular they are not variations of the writer marginal.
For equal-mass finite measures, use
$\pTV(\nu,\nu')=\frac12|\nu-\nu'|(\mathcal Q)$; this also covers
subprobabilities of equal mass.

Write $V_i(f)=\int\pVar_i(f)$ for the directional BV seminorm, with
spectators and $X$ integrated. The same definitions apply to a restricted
subdensity of total mass $m\leq1$. Only the explicitly active directions
are varied; jumps in a fixed old archive direction are not included.

\begin{lemma}[Grid and sectional concentration]\label{p1p:lem:grid}
Let a nonnegative subdensity $h$ have integrated mass $m$ and active
two-coordinate variation $V=V_1+V_2$. Let $P_N$ average on the $N$ by $N$
grid at fixed spectators. Then
\begin{equation}\label{p1p:eq:grid}
 \|h-P_Nh\|_1\leq V/(2N),\qquad V_i(P_Nh)\leq V_i(h).
\end{equation}
On each active unit-square fibre, with mass $m_z$ and variations
$V_{1,z},V_{2,z}$,
\begin{equation}\label{p1p:eq:sectionL2}
 \|h_z\|_2\leq\sqrt{(m_z+V_{1,z})(m_z+V_{2,z})}
 \leq m_z+(V_{1,z}+V_{2,z})/2.
\end{equation}
If $K$ preserves area on each such fibre, fixes all spectators, and has
active displacement $\leq\delta$ off a set of sectional area $\leq\beta$,
then
\begin{equation}\label{p1p:eq:capfreegrid}
 \pTV(K_\#(h\lambda),h\lambda)
 \leq \frac{V}{2N}+(m+V/2)\sqrt{\beta+4N\delta}.
\end{equation}
If instead $h\leq C$, the alternative bound is
\begin{equation}\label{p1p:eq:capgrid}
 \pTV(K_\#(h\lambda),h\lambda)
 \leq V/(2N)+2CN\delta+C\beta/2.
\end{equation}
\end{lemma}
\begin{proof}
For an interval $I$ of length $l$ and its mean $h_I$,
\[
 \int_I|h-h_I|\leq l^{-1}\int_{I\times I}|h(x)-h(y)|\,dx\,dy
 \leq (l/2)|Dh|(I).
\]
The last inequality follows by integrating the derivative between $x$
and $y$: its weight at $t\in I$ is
$2(t-\inf I)(\sup I-t)/l\leq l/2$.
Successive coordinate averaging contracts $L^1$ and the unused
directional variations, proving the first bound in~\eqref{p1p:eq:grid}.
For adjacent cell means the difference is the average of
$h(t+1/N)-h(t)$. Integrating its derivative along that segment and summing
the cells gives overlap weight at most one. This proves the variation
contraction, first for smooth functions and then for BV by approximation.

For a BV square slice the one-dimensional representatives give, almost
everywhere,
\[
 h(x,y)\leq A(y):=\int_0^1h(t,y)dt+\pVar_x h(\cdot,y),
 \quad
 h(x,y)\leq B(x):=\int_0^1h(x,t)dt+\pVar_y h(x,\cdot).
\]
The common representatives exist by BV slicing. Thus $h^2\leq A(y)B(x)$;
Fubini proves the first inequality in~\eqref{p1p:eq:sectionL2}, and the
arithmetic--geometric mean proves the second.

Put $g=P_Nh$. Its value can change under $K$ only on the bad set or when
an active coordinate crosses a grid line. This exceptional set $E$ has
sectional area at most $s=\beta+4N\delta$. Since $K$ preserves sectional
area, so does its image. Cauchy bounds each of
$\int_Eg$ and $\int_Eg\circ K$ by
$\sqrt{s}\|g\|_2$. Their sum is divided by two in total variation.
The two approximation errors together are
$\|h-g\|_1\leq V/(2N)$, giving~\eqref{p1p:eq:capfreegrid} after integrating
the section coefficients. If $0\leq g\leq C$, then
$|g-g\circ K|\leq C$ on $E$, yielding~\eqref{p1p:eq:capgrid} instead.
The argument prices both the exceptional set and its image.
\end{proof}

Let $\mathcal T_n$ be the complete composition of the exact cycles and
$\mathcal B_n$ the corresponding baker composition. Let $R_n$ be the
final ready quantile and $Z_n=(X,B_1,\ldots,B_n,W)$ all retained external
and positional nuisance variables. The internal memories are retained
in the quantum wave; there is no additional hidden internal coordinate
to marginalize in the specified ontology.

\begin{theorem}[Quantitative preparation with retained archives]
\label{p1p:thm:prep}
Restrict the actual entrance law to a subdensity $h_x\leq f_x$ of mass
$m=1-\tau$, with directional variations $V_w,V_1,\ldots,V_n$ in the writer
and warmup archives. Set
\begin{equation}\label{p1p:eq:Vsum}
 S_n=2V_w+\sum_{j=1}^n V_j,\qquad H_n^*=n+S_n/2,
 \qquad E_n=\frac{S_n}{2N}+H_n^*\sqrt{\beta+4N\delta}.
\end{equation}
Suppose the identical exact cycle has the sectional comparison bounds
of Lemma~\ref{p1p:lem:separation}. Then
\begin{align}
 \pTV\bigl(\pLaw(R_n,Z_n),\pUnif\otimes\pLaw(Z_n)\bigr)
 &\leq\tau+\frac{V_w}{4\,2^n}+2E_n,
 \label{p1p:eq:freshness}\\
 \Pr(\hbox{some archive miscopies its own writer sign})
 &\leq\tau+H_n^*\sqrt\beta+nE_n.
 \label{p1p:eq:allrecords}
\end{align}
Each correctly copied archive sign persists throughout all later
prescribed stages and holding intervals. If the full entrance density
is capped by $C$, one may instead take $\tau=0$ and
\begin{equation}\label{p1p:eq:cappedE}
 E_n^{\rm cap}=\frac{S_n}{2N}+2nCN\delta+nC\beta/2.
\end{equation}
In that capped case the all-record failure also has the simpler bound
$nC\beta$.
\end{theorem}
\begin{proof}
After $j$ ideal cycles, the archive values determine
$s_i=\lfloor2b_i\rfloor$ and $v_i=2b_i-s_i$ for $1\leq i\leq j$.
Writing $H_j=\sum_{i=1}^j2^{j-i}s_i$, the inverse writer coordinate is
$u_0=2^{-j}(u+H_j)$. The full inverse Jacobian is one: the writer factor
$2^{-j}$ cancels the $j$ archive factors two. At fixed old archives,
differentiation in the active writer therefore multiplies its variation
by $2^{-j}$, whereas variation in the next unused $v_{j+1}$ is unchanged.
Any old archive-digit seam is in a fixed coordinate. Thus the sum of
active variations before all $n$ ideal steps is at most $S_n$.

Telescope $\mathcal T_n-\mathcal B_n$ by replacing one ideal step at a
time, using ideal prefixes and exact suffixes. Deterministic pushforward
contracts total variation, so Lemma~\ref{p1p:lem:grid} bounds the total
boxed endpoint discrepancy by $E_n$; using $nm\leq n$ gives
\eqref{p1p:eq:Vsum}. The same bound controls every partial-prefix error.

For the complete ideal output, fix all original source values and $X$.
The archives reveal the dyadic cell containing $u_0$, as well as those
source values. Replacing the original density by its cell mean in $u_0$
is exactly the operation that makes the remainder uniform while keeping
the ideal archive marginal. The interval inequality in the proof of
Lemma~\ref{p1p:lem:grid} gives $L^1$ error at most $2^{-n}V_w/2$ and hence
TV error at most $V_w/(4\,2^n)$. Transferring both the complete endpoint
law and its archive marginal from ideal to exact costs $2E_n$.

The removed entrance subprobability and the product of its own final
nuisance marginal with $\pUnif$ each have mass $\tau$. Their distance is
at most $\tau$. This proves~\eqref{p1p:eq:freshness} with that tail charged
once, not once per cycle.

For each record, the wrong-copy event is contained in its pre-copy bad
set. At the ideal boxed prefix, sectional Cauchy bounds its mass by
$(m+V_{{\rm active},j}/2)\sqrt\beta$. Replacing this prefix by the actual
boxed prefix costs at most $E_n$. Sum over $j$, add the initial tail once,
and obtain~\eqref{p1p:eq:allrecords}. The zero archive current established
above proves subsequent persistence. Under a cap, exact area preservation
keeps that same cap at every actual prefix, giving $nC\beta$ directly.
The capped grid lemma gives~\eqref{p1p:eq:cappedE}.
\end{proof}

If the averaged complete quantile Fisher information is
$J^{\rm q}=\int\sum_i|\partial_{u_i}\log f_x|^2 f_x\,d\lambda\,d\mu_X$,
score Cauchy gives
\begin{equation}\label{p1p:eq:quantilescore}
 S_n\leq\sqrt{n+4}\sqrt{J^{\rm q}},\qquad V_w\leq\sqrt{J_w^{\rm q}}.
\end{equation}
This is one sufficient regularity condition. The physical-score result
below is strictly broader and does not assume finite quantile Fisher.

\section{Physical relative scores without a density cap}
\label{p1p:sec:physical}

Let $\xi=(\xi_0,\ldots,\xi_{d-1})$ be all standardized physical entrance
coordinates and let $\pGam_d(\xi)=\prod_i\phi(\xi_i)$. Conditional on the
same retained $X=x$, write the actual density as
\begin{equation}\label{p1p:eq:relative}
 p_x(\xi)=\pGam_d(\xi)f_x(\xi),\qquad
 \partial_i f_x=f_xs_{i,x}\quad\hbox{distributionally},\qquad
 J_i=\int |s_{i,x}|^2p_x\,d\xi\,d\mu_X(x)<\infty.
\end{equation}
The score is assigned arbitrarily on zero-density sets. In particular
$\sum_iJ_i$ concerns the full conditional density, including every
archive and unused future source. The comparison Gaussian fixes the
analysis coordinates and describes the quantum stock; it does not
specify the actual law $p_x$.

\begin{lemma}[Moments and boundary traces from weak scores]
\label{p1p:lem:physicalBV}
Under~\eqref{p1p:eq:relative}, the actual second moment $M_i=\mathbb E\xi_i^2$
is finite and the ordinary physical Fisher information satisfies
\begin{equation}\label{p1p:eq:momentidentity}
 I_i^{\rm abs}=J_i+2-M_i\geq0,
 \qquad M_i\leq J_i+2.
\end{equation}
Restrict the complete actual density to the physical box
$|\xi_i|\leq R_i$ and extend by zero in the quantile coordinates
$u_i=\Phi(\xi_i)$. Its directional variations obey
\begin{equation}\label{p1p:eq:physicalBV}
 V_i\leq D_{R_i}\{\sqrt{J_i}+\sqrt{J_i+2}\},
 \qquad D_R=\sqrt{2\pi}e^{R^2/2}.
\end{equation}
For a common radius $R$, total score $J=\sum_iJ_i$, one writer and $n$
active archives,
\begin{equation}\label{p1p:eq:physicalS}
 S_n\leq D_R\sqrt{n+4}\{\sqrt J+\sqrt{J+2d}\}.
\end{equation}
\end{lemma}
\begin{proof}
First prove finiteness rather than assuming the integration by parts is
legitimate. Choose even compact smooth cutoffs $0\leq\chi_R\leq1$,
nonincreasing in $|\xi_i|$, increasing to one, with uniformly bounded
$\xi_i\chi_R'$. Set $M_R=\mathbb E(\xi_i^2\chi_R)$.
Integration by parts against the compact test $\xi_i\chi_R$ gives
\[
 M_R=\mathbb E(\chi_R+\xi_i\chi_R')+
       \mathbb E(\xi_i\chi_Rs_i)
 \leq1+\sqrt{M_RJ_i}.
\]
Spectator cutoffs can be exhausted: the active test is bounded and
its score term is integrable by Cauchy. Solving the quadratic gives
$M_R\leq[(\sqrt{J_i}+\sqrt{J_i+4})/2]^2$; monotone convergence proves
$M_i<\infty$. Now $\xi_i s_i$ is integrable. Removing the cutoff, its
bounded derivative term tends to zero by dominated convergence, giving
$\mathbb E\xi_i s_i=M_i-1$. The ordinary score is $s_i-\xi_i$, so
\[
 I_i^{\rm abs}=\mathbb E(s_i-\xi_i)^2=J_i+2-M_i.
\]
This proves~\eqref{p1p:eq:momentidentity} without presupposing the moment.

For almost every $x$, the one-coordinate marginal $p_{i,x}$ has weak
derivative equal to the integrated joint derivative. Conditional
expectation and Cauchy imply
$\|p_{i,x}'\|_1\leq\sqrt{I_{i,x}^{\rm abs}}$.
An integrable nonnegative $W^{1,1}(\mathbb R)$ density tends to zero at
both infinities, so
$2\sup p_{i,x}\leq\|p_{i,x}'\|_1$.
On the box interior, differentiating in $u_i$ introduces
$1/\phi(\xi_i)\leq D_{R_i}$; score Cauchy bounds the integrated interior
variation by $D_{R_i}\sqrt{J_i}$.
At each of the two quantile faces the zero extension contributes the
trace of the density. Integrating all other coordinates bounds their
sum by
\[
 D_{R_i}\int[p_{i,x}(R_i)+p_{i,x}(-R_i)]\,d\mu_X(x)
 \leq D_{R_i}\sqrt{J_i+2}.
\]
BV traces or a limiting regular face justify the same assertion for weak
densities. This proves~\eqref{p1p:eq:physicalBV}. Finally apply weighted
Cauchy to the coefficients $(2,1,\ldots,1)$ in $S_n$. Their squared sum
is $n+4$; including unused coordinates in $J$ and $d$ only enlarges the
upper bound. This proves~\eqref{p1p:eq:physicalS}.
\end{proof}

Higher moments, when assumed, give a useful separate tail estimate:
\begin{equation}\label{p1p:eq:tail}
 \tau\leq\sum_i\frac{\mathbb E|\xi_i|^p}{R_i^p}.
\end{equation}
The second-moment part of this estimate follows from the weak-score
hypothesis. A twentieth moment does not follow from that argument and
will be explicitly required for the next row.

\begin{corollary}[A finite cap-free preparation row]\label{p1p:cor:row}
Include one writer, $100$ warmup archives, and one untouched subsequent
measurement source, so $d=102$. Suppose the actual complete conditional
law satisfies
\[
 \sum_iJ_i\leq100,\qquad \sum_i\mathbb E|\xi_i|^{20}\leq10^{15}.
\]
Use $D=A=50$, the complete entrance box $|\xi_i|\leq10$, and analysis
grid $N=10^{40}$. Then the preparation and all later current-zero holds
satisfy
\begin{align}
 E_{100}&\leq2.00006\times10^{-16},\label{p1p:eq:rowE}\\
 \Delta_{\rm box}:=\frac{V_w}{4\,2^{100}}+2E_{100}
 &<5.192874163845\times10^{-8},\label{p1p:eq:rowbox}\\
 \pTV(\pLaw(R_{100},Z_{100}),\pUnif\otimes\pLaw(Z_{100}))
 &<1.005192874164\times10^{-5},\label{p1p:eq:rowfresh}\\
 \Pr(\hbox{some wrong warmup copy})
 &<1.000000002000061\times10^{-5}.\label{p1p:eq:rowcopy}
\end{align}
The untouched measurement source is part of $Z_{100}$, not averaged out.
\end{corollary}
\begin{proof}
The elementary exponential enclosure gives $D_{10}<1.31\times10^{22}$.
Using $\sqrt{102}<10.1$, $\sqrt{104}<10.2$ and $\sqrt{304}<17.44$ in
Lemma~\ref{p1p:lem:physicalBV} gives
\[
 V_w<2.6331\times10^{23},\quad
 S_{100}<3.6665328\times10^{24}<3.67\times10^{24},\quad
 H_{100}^*<2\times10^{24}.
\]
The last inequality retains the positive mass term $100$ in $H_{100}^*$.
Choose $c=3$, $\eta=10^{-110}$, $k=24$. Gaussian upper tails and
exponential bounds give
\begin{align*}
 \gamma=e^{-300}&<10^{-130},&
 \Phi(-24)&<\eta/2,\\
 \rho_A=e^{-2600}&<10^{-1100},&
 t_A=\Phi(-76)&<10^{-1200},\\
 a_D=\Phi(-53)&<10^{-600}.&&
\end{align*}
Thus $\delta<2\times10^{-130}$ and
$\beta<3\times10^{-110}$. In particular
$\sqrt{\beta+4N\delta}<3\times10^{-45}$ and
$\sqrt\beta<2\times10^{-55}$. Therefore
\[
 E_{100}\leq\frac{4\times10^{24}}{2\times10^{40}}
       +(2\times10^{24})(3\times10^{-45})
       =2.00006\times10^{-16}.
\]
The ideal boxed writer fee is bounded by
$2.6331\times10^{23}/(4\,2^{100})$
$<5.192874123844\times10^{-8}$.
The actual complete-box tail is at most $10^{15}/10^{20}=10^{-5}$.
Substitution in Theorem~\ref{p1p:thm:prep} gives the stated rational
decimal ceilings; the copy ceiling includes the additional
$4\times10^{-31}$ bad-set term.

For reproducible directed inequalities one may use
$\sqrt{2\pi}<2.51$ and bound $e^{50}$ by its positive Taylor sum through
degree $200$ plus the next term divided by $1-50/202$.
For the lower exponential bounds used in the tails,
$e>\sum_{j=0}^5 1/j!=163/60$ suffices, together with
$\Phi(-x)\leq e^{-x^2/2}/(2x)$ for $x>0$.
All row decisions consequently reduce to rational inequalities.
\end{proof}

An explicit non-Born member of this class consists of $102$ independent
actual centered Gaussians of variance $5/2$ in the standardized physical
coordinates. Its full relative Fisher is
\[
 J=102(1-2/5)^2(5/2)=91.8,
\]
and its twentieth-moment sum is
\[
 102\,(19!!)\,(5/2)^{10}=6.368862690925598\ldots\times10^{14}<10^{15}.
\]
Its ground-relative density is unbounded. Its quantile Fisher is
infinite: for one coordinate the integrand contains
$x^2p(x)/\phi(x)^2$, whose exponential factor grows at infinity.
Independence is used only to exhibit this member; neither the theorem
nor its constants impose it.

\subsection{Finite-resource existence and growing conditional prefixes}

\begin{theorem}[Uniform writer-score family and all retained records]
\label{p1p:thm:existence}
Let a consistent stock family provide, for every finite $n$, the full
conditional density of the writer, $n$ warmup sources, one unused future
source, and retained $X$. Suppose its weak relative scores satisfy
\[
 J_w(n)\leq J_*<\infty\quad\hbox{for every }n,
 \qquad J_i(n)<\infty\quad\hbox{for every other coordinate of each prefix}.
\]
Then for every $0<\varepsilon<1$ there are finite $n$, finite physical
cutoffs and finite separations $D,A$ for which complete retained-archive
freshness is at most $3\varepsilon/4$ and the probability of any wrong
warmup copy-and-hold record is less than $0.42\varepsilon$.
\end{theorem}
\begin{proof}
Choose $R_w^2\geq8(J_*+2)/\varepsilon$. The writer tail is at most
$\varepsilon/8$ and
\[
 V_w\leq V_*:=D_{R_w}(\sqrt{J_*}+\sqrt{J_*+2}),
\]
independently of $n$. Choose finite $n\geq1$ so that
$V_*/(4\,2^n)\leq\varepsilon/4$. For this chosen full prefix, choose
each of the $n+1$ source cutoffs to satisfy
$R_i^2\geq8(n+1)(J_i(n)+2)/\varepsilon$.
Their \emph{aggregate} tail is at most $\varepsilon/8$; hence
$\tau\leq\varepsilon/4$. All boxed source variations are finite by
Lemma~\ref{p1p:lem:physicalBV}. Set $S=S_n$ and $H=n+S/2$ and choose
\begin{equation}\label{p1p:eq:schedule}
 N\geq\max\{1,8nS/\varepsilon\},\qquad
 \beta\leq\frac{\varepsilon^2}{512n^2H^2},\qquad
 \delta\leq\frac{\varepsilon^2}{2048n^2NH^2}.
\end{equation}
The two terms in $E_n$ are at most $\varepsilon/(16n)$ each, so
$E_n\leq\varepsilon/(8n)$. Thus freshness is at most
$\varepsilon/4+\varepsilon/4+\varepsilon/(4n)\leq3\varepsilon/4$,
and the all-record bound is at most
\[
 \varepsilon\left(\frac14+\frac{1}{\sqrt{512}\,n}+\frac18\right)
 \leq\varepsilon\left(\frac38+\frac4{89}\right)
 =\frac{299}{712}\varepsilon<0.42\varepsilon.
\]
Here $\sqrt{512}>89/4$. The factor $n$ in~\eqref{p1p:eq:schedule}
prices all own-event records; a schedule controlling freshness alone
does not automatically do so.

It remains to realize these positive $\beta,\delta$ thresholds with
finite Gaussian parameters. Fix $c=3$ and choose $\eta$ sufficiently
small for the source-end part of $\beta$. Choose finite $k$ with
$\Phi(-k)<\eta/2$. Increase $D$ until $e^{-2Dc}$ is below both the
desired displacement allocation and $\eta/2$, and until
$\Phi(-(D-c))$ and $\Phi(-(D+c))$ meet their respective allocations.
Increase $A>k$ until $e^{-2A(A-k)}$ and $\Phi(-(2A-k))$ meet the remaining
allocations. Lemma~\ref{p1p:lem:separation} proves the required map and
copy bounds. All choices are finite after this finite prefix has been
selected. The translations and conjugated holding swaps described
above implement the corresponding smooth prescribed parent.
\end{proof}

The order of these choices matters. A common second-moment box whose
total moment grows like $b(n+1)$ would require
$R^2\geq b(n+1)/\tau$. Its writer variation enclosure contains
$e^{b(n+1)/(2\tau)}$, which can grow faster than $2^n$. The theorem fixes
the writer cutoff before increasing the source count, assigning costly
source variations to the later grid and separation choices.

\begin{proposition}[Why finite-prefix smoothness is insufficient]
\label{p1p:prop:prefixcounter}
There is a consistent actual stock family for which every finite
physical prefix has a smooth positive density, finite relative Fisher
and bounded coordinate second moments, but the complete ideal baker
output satisfies, for every $n\geq1$,
\[
 \pTV(\pLaw(R_n,Z_n),\pUnif\otimes\pLaw(Z_n))>
 \frac{5189}{12550}>0.4.
\]
\end{proposition}
\begin{proof}
Take independent standard normals $Y,G_1,G_2,\ldots$ and set
$Z_j^{\rm in}=Y+2^{-6j}G_j$; use writer $u=\Phi(Y)$ and sources
$v_j=\Phi(Z_j^{\rm in})$. Every finite covariance is nonsingular.
For a prefix containing the writer and $m$ sources, differentiation at
fixed source coordinates gives
\[
 \partial_Y\log(p/\pGam)=\sum_{j=1}^m2^{12j}(Z_j^{\rm in}-Y),
 \qquad J_w^{(m)}=\sum_{j=1}^m2^{12j}.
\]
A bank with $n$ warmup archives and an untouched receiver has $m=n+1$.
Thus this family fails the uniform full conditional writer hypothesis,
although its writer marginal is exactly standard normal.
The final archive $b_n$ reveals $v_n=2b_n-\lfloor2b_n\rfloor$, hence
$Z_n^{\rm in}$. Define the archive prediction
$\widehat r=\{2^n\Phi(Z_n^{\rm in})\}$.
Circle distance obeys
\[
 \operatorname{dist}_{\mathbb T}(r,\widehat r)
 \leq 2^{-5n}|G_n|/\sqrt{2\pi}.
\]
Since $\sqrt{2\pi}>2.5$, the event $|G_n|\leq0.8$ implies this distance
is at most $0.01$ for all $n\geq1$. Under an independent uniform
remainder, the same archive-dependent circle interval has probability
$0.02$. Meanwhile
\[
 \Pr(|G_n|\leq0.8)
 >\frac{1.6(1-0.32)}{2.51}.
\]
Subtracting $0.02$ gives $5189/12550$. The strict inequality follows
from $e^{-0.32}>1-0.32$ and $\sqrt{2\pi}<2.51$.
\end{proof}

\subsection{Two distinct cap-free limits}

\begin{proposition}[Fixed-law finite-map convergence]\label{p1p:prop:L1}
Fix $n$ and any complete entrance law with conditional $L^1$ densities
$f_x$. As $D,A\to\infty$,
\[
 \pTV((\mathcal T_n)_\#\mu,(\mathcal B_n)_\#\mu)\longrightarrow0.
\]
No cap or Fisher condition is needed for this fixed-law assertion.
\end{proposition}
\begin{proof}
For fixed $u<1/2$ and $v\in(0,1)$,
$F_D^{-1}(u)=-D+\Phi^{-1}(2u)+o(1)$ and its minority posterior tends to
zero. The copy position is $-A+\Phi^{-1}(v)+o(1)$, and the return and
archive outputs tend to $(2u,v/2)$. Reflection gives the other branch.
Consequently $B^{-1}T_{D,A}\to I$ almost everywhere. Finite iteration
excludes only the finite family of dyadic cuts and their preimages.
The corresponding full difference maps preserve Lebesgue measure.
For bounded continuous $h$ on the closed cube, dominated convergence
gives $\|h\circ K-h\|_1\to0$. Approximate an arbitrary $f\in L^1$ by such
$h$ and use
\[
 \|f\circ K-f\|_1\leq2\|f-h\|_1+\|h\circ K-h\|_1.
\]
This proves TV convergence on each $X$ fibre. Dominated convergence in
the actual $X$ law completes the argument.
\end{proof}

A qualitative growing-family sufficient condition is stronger: require
that $U$ conditional on the \emph{entire} consistent source tape and $X$
has an $L^1$ density. Dyadic cell averaging converges in $L^1$ on each
such fibre. Its error is precisely the ideal complete-remainder
freshness error before projecting to a finite retained tape. Dominated
convergence allows a finite $n$ to be selected for each tolerance;
Proposition~\ref{p1p:prop:L1} then selects finite separations for that
fixed prefix. Proposition~\ref{p1p:prop:prefixcounter} explains why
absolute continuity of every finite prefix is weaker than this premise.

\begin{proposition}[A finite-resource obstruction without a density cap]
\label{p1p:prop:nocapcounter}
For $n=44$ and $D=A=12$, there is a smooth actual complete law with
physical relative Fisher exactly $100$ whose final marginal ready
distance from uniform exceeds $0.94$.
\end{proposition}
\begin{proof}
Take actual writer $Q_0\sim N(10,1)$ and independent standard-normal
actual archives. The entrance ratio is $e^{10Q_0-50}$, with physical
relative score $10$ in the writer direction and zero in every archive
direction. Its relative density is unbounded.
The Gaussian upper-tail inequality implies
$\Phi(-8.1)<0.005/2^{44}$ and
$\Pr(Q_0>8.1)>0.96$. On this event every ideal digit is right and every
ideal remainder is greater than $0.995$. For each actual archive,
exclude its initial quantiles outside $[10^{-12},1-10^{-12}]$; the
aggregate exclusion probability is at most $88\times10^{-12}$.
With $c=3,k=8$, Lemma~\ref{p1p:lem:separation} gives, on the right good
branch, a ready-coordinate error per cycle at most
$10^{-48}+10^{-41}$. Induction bounds the accumulated error by
$(2^{44}-1)(10^{-48}+10^{-41})<10^{-20}$.
Every actual remainder therefore stays in the right good branch and
the final one exceeds $0.99$. Thus
\[
 \Pr(R_{44}>0.99)>0.96-88\times10^{-12},
\]
whereas uniform assigns $0.01$ to this event. The difference exceeds
$0.94$. Complete archive-conditioned distance is at least this marginal
distance by projection. This is a finite-resource counterexample, not
a failure of an upper estimate.
\end{proof}

For comparison, the cap-based physical-score row does survive with
$C=10$, $J\leq100$, $R=6$, $n=44$, $D=A=12$ and $N=10^{19}$.
The direct cap trace bound is
$V_i\leq D_R\sqrt{J_i}+2C$, and the complete-box tail is
$2Cd\Phi(-R)$. With $d=46$, including the untouched future source,
$D_6<1.65\times10^8$ and $\Phi(-6)<10^{-9}$ give
$S_{44}<1.2\times10^{10}$ and $V_w<1.65000002\times10^9$.
Using $\delta<2\times10^{-30}$,
$\beta\leq2\times10^{-12}+10^{-18}$ in~\eqref{p1p:eq:cappedE} gives
$E_{44}^{\rm cap}\leq1.864000022\times10^{-8}$ and freshness
$<2.440519056475\times10^{-5}$.
Alternatively, the complete quantile-score row $n=18,C=10,J^{\rm q}\leq10^4$
with $N=10^{15}$ gives $E_{18}^{\rm cap}\leq1.8097009\times10^{-10}$ and
freshness $\leq9.5367793580805\times10^{-5}$.
These are different admitted-law classes; the cap-free row of
Corollary~\ref{p1p:cor:row} changes both the class and the resources.

\section{Robustness under complete physical map errors}
\label{p1p:sec:robustness}

The exact Gaussian parent supplies the complete map of
Theorem~\ref{p1p:thm:map}. A replacement parent must control its complete
configuration transport, including sources, clocks and any archive
coordinates that cease to be stationary. The following results state
what such a control implies for the admitted actual laws.

\begin{lemma}[Reference error and actual density]\label{p1p:lem:QV}
Let $K$ be a measurable map of the unit cube that fixes all but $k$
coordinates. Suppose
\[
 \pTV(K_\#\lambda,\lambda)\leq e_{\rm ref},
 \qquad |K_i(u)-u_i|\leq\delta
\]
in each active direction outside a set of reference volume $\beta$.
For a nonnegative subdensity $h\leq C$ with active variation $V$,
\begin{equation}\label{p1p:eq:QV}
 \pTV(K_\#(h\lambda),h\lambda)
 \leq\frac{V}{2N}+C\{e_{\rm ref}+kN\delta+\beta/2\}.
\end{equation}
No bijectivity or reference preservation of $K$ is required.
\end{lemma}
\begin{proof}
Let $g=P_Nh$ use the $N$-cell grid in the active directions.
The approximation costs before and after pushforward total at most
$V/(2N)$. There is an exact signed-measure decomposition
\begin{equation}\label{p1p:eq:decomposition}
 K_\#(g\lambda)-g\lambda
 =K_\#[(g-g\circ K)\lambda]+g(K_\#\lambda-\lambda).
\end{equation}
Here the last multiplication is at the output coordinate. The first
half-variation is at most $\frac12\int|g-g\circ K|d\lambda$.
The integrand vanishes except on the bad set and grid-crossing strips,
of total volume at most $\beta+2kN\delta$. It is at most $C$ there.
The second half-variation is at most $Ce_{\rm ref}$. This proves
\eqref{p1p:eq:QV} even for merging maps.
\end{proof}

For two conservative guidance endpoint maps $\mathcal F_0,\mathcal F$
on the same complete configuration space and with the same entrance
reference, take $K=\mathcal F_0^{-1}\mathcal F$ in entrance reference
coordinates. The ideal inverse must be defined almost everywhere on
the physical output, including any leakage. If it is a common-space
bijection on that support, equivariance and invariance of TV under
bijective coordinates give
\[
 e_{\rm ref}=\pTV((\mathcal F)_\#\lambda,(\mathcal F_0)_\#\lambda)
 =\pTV(\rho_{\rm phys},\rho_0).
\]
For normalized complete endpoint waves,
\begin{equation}\label{p1p:eq:waveReference}
 e_{\rm ref}\leq\|\Psi-\Psi_0\|_2.
\end{equation}
Indeed $||\Psi|^2-|\Psi_0|^2|\leq
(|\Psi|+|\Psi_0|)|\Psi-\Psi_0|$ and Cauchy give the bound after division
by two. This controls reference mass transport. The displacement and
exceptional-set inputs to Lemma~\ref{p1p:lem:QV} still concern the true
complete maps and require their own estimates.

\subsection{Full-dimensional concentration from physical scores}

\begin{lemma}[A concentration coefficient for arbitrary external memory]
\label{p1p:lem:concentration}
Let $d\geq2$, $p=d/(d-1)$, and let $h_x$ be the boxed conditional actual
density in physical reference quantiles. Write its mass and variations
as $m_x\leq1$ and $V_{i,x}$, and define
\begin{equation}\label{p1p:eq:Bcoefficient}
 B_x=m_x+\frac1d\sum_i V_{i,x}.
\end{equation}
Then
\begin{equation}\label{p1p:eq:Lp}
 \|h_x\|_{L^p(\lambda)}\leq
 \prod_i(m_x+V_{i,x})^{1/d}\leq B_x.
\end{equation}
Under a common physical cutoff $R$ and the averaged complete score bound
$\sum_iJ_i\leq J$,
\begin{equation}\label{p1p:eq:B2}
 \|B_x\|_{L^2(\mu_X)}\leq
 B_2:=1+\frac{D_R}{\sqrt d}\{\sqrt J+\sqrt{J+2d}\}.
\end{equation}
\end{lemma}
\begin{proof}
For a nonnegative BV cube function define
\[
 A_i(u_{-i})=\int_0^1h(u)du_i+\pVar_i h(\cdot,u_{-i}).
\]
The slice inequality gives $h\leq A_i$ simultaneously almost everywhere.
The product integration inequality
\begin{equation}\label{p1p:eq:productHolder}
 \int\prod_{i=1}^d A_i(u_{-i})^{1/(d-1)}du
 \leq\prod_{i=1}^d\left(\int A_i\right)^{1/(d-1)}
\end{equation}
follows by induction. For $d=2$ it is Fubini. For $d>2$, integrate $u_d$
first and apply H\"older to the $d-1$ factors depending on it. Denote their
integrals by $F_i$, $i<d$. On the remaining coordinates separate
$A_d^{1/(d-1)}$ by H\"older with exponents $d-1$ and
$(d-1)/(d-2)$, and apply the dimension-$d-1$ inequality to the $F_i$.
This yields precisely~\eqref{p1p:eq:productHolder}.
Multiplying the bounds $h\leq A_i$, integrating, and raising to power
$(d-1)/d$ proves the first part of~\eqref{p1p:eq:Lp}. The second is the
arithmetic--geometric mean. BV approximation preserves these bounds.

Conditionally on $X=x$, Lemma~\ref{p1p:lem:physicalBV} gives
\[
 \sum_iV_{i,x}\leq D_R\sqrt d
 \{\sqrt{J_x}+\sqrt{J_x+2d}\},\qquad J_x=\sum_iJ_{i,x}.
\]
Apply Minkowski in $L^2(\mu_X)$, use $m_x\leq1$ and
$\int J_xd\mu_X\leq J$. This proves~\eqref{p1p:eq:B2}.
Thus the square-integrable coefficient is obtained from the averaged
physical score; it is not inferred from an averaged BV bound alone.
\end{proof}

\begin{theorem}[Cap-free robustness from reference TV]
\label{p1p:thm:CFQ}
On each external-memory fibre let $K_x$ satisfy
Lemma~\ref{p1p:lem:QV} with errors $e_{{\rm ref},x},\delta_x,\beta_x$.
Set
\[
 r_x=e_{{\rm ref},x}+kN\delta_x+\beta_x/2,
 \qquad\overline r=\int r_x\,d\mu_X(x).
\]
For the boxed actual law, with integrated active variation $V$,
\begin{equation}\label{p1p:eq:CFQ}
 \ell_{\rm box}:=\pTV(K_\#(h\lambda\mu_X),h\lambda\mu_X)
 \leq\frac{V}{2N}+2B_2\overline r^{\,1/d}.
\end{equation}
All conditional error averages here use the actual external-memory law.
\end{theorem}
\begin{proof}
For a proof threshold $M>0$ truncate $h_x$ to $h_{x,M}=\min(h_x,M)$.
This contracts each BV seminorm. The discarded mass is at most
\[
 \tau_{x,M}\leq M^{-(p-1)}\int h_x^p
 \leq B_x^pM^{-(p-1)}.
\]
The discarded input and its pushforward have the same mass, so their
distance costs at most $\tau_{x,M}$ once. Lemma~\ref{p1p:lem:QV} gives
\[
 \ell_x\leq\frac{V_x}{2N}+Mr_x+B_x^pM^{-(p-1)}.
\]
Choose $M=B_xr_x^{-(d-1)/d}$; the two last terms are each
$B_xr_x^{1/d}$. If $r_x=0$, let $M\to\infty$; if $B_x=0$ there is no mass.
Average in $x$ and apply Cauchy and concavity of $t^{2/d}$:
\[
 \int B_xr_x^{1/d}d\mu_X
 \leq\|B_x\|_2\left(\int r_x^{2/d}d\mu_X\right)^{1/2}
 \leq B_2\overline r^{1/d}.
\]
This proves the theorem without imposing a cap on the actual density.
\end{proof}

For a common boxed ideal baseline with remainder freshness
$\Delta_{0,\rm box}$, the physical law has own-marginal freshness at most
\begin{equation}\label{p1p:eq:physicalfresh}
 \tau+\Delta_{0,\rm box}+2\ell_{\rm box}.
\end{equation}
As before, one law error moves the joint distribution and the other
moves its actual nuisance marginal. The complete initial tail is
compared directly with its own-marginal product and is charged once.
For the $d=102$ row, $V<4\times10^{24}$ and $B_2<4\times10^{22}$.
For example $N=10^{32}$, $k\leq102$, and actual-memory averaged errors
\[
 \overline e_{\rm ref}\leq10^{-3200},\quad
 \overline\beta\leq10^{-3200},\quad
 \overline\delta\leq10^{-3300}
\]
give $\overline r<10^{-3162}=(10^{-31})^{102}$ and
$\ell_{\rm box}\leq2.8\times10^{-8}$. Equation~\eqref{p1p:eq:physicalfresh}
then yields a freshness ceiling $1.0108\times10^{-5}$.
This row quantifies the severe dimension loss in an ordinary
reference-TV criterion. It specifies a mathematical tolerance interface;
it does not assert that a finite material source achieves those errors.

\subsection{A sharper two-coordinate distortion criterion}

\begin{theorem}[Sectional density calibration]\label{p1p:thm:sectiondist}
Suppose $K$ fixes all but two coordinates, and on each spectator fibre
$z$ the pushforward of active reference area has density $r_z$ with
\[
 \epsilon_{2,z}=\|r_z-1\|_{L^2((0,1)^2)}<\infty.
\]
Suppose also that the active displacement is at most $\delta_z$ off a
set of sectional area at most $\beta_z$. Define
$H_z=m_z+(V_{1,z}+V_{2,z})/2$. Then
\begin{equation}\label{p1p:eq:sectiondist}
 \ell_{\rm box}\leq\frac{V}{2N}
 +\int H_z\{\sqrt{\beta_z+4N\delta_z}+\epsilon_{2,z}\}\,d\zeta(z),
\end{equation}
where $d\zeta$ is reference measure in the fixed physical spectators
times the actual external-memory law. Under uniform sectional errors,
\begin{equation}\label{p1p:eq:sectionuniform}
 \ell_{\rm box}\leq\frac{V}{2N}
 +(m+V/2)\{\sqrt{\beta+4N\delta}+\epsilon_2\}.
\end{equation}
\end{theorem}
\begin{proof}
On a fixed fibre put $g=P_Nh$ and let $E$ be its bad/crossing set, of
area at most $s=\beta+4N\delta$. The subreference pushforward
$K_\#(1_E\lambda)$ has a density $a\leq r$ and mass at most $s$.
Decompose $a=\min(a,1)+(a-1)_+$. Since
$\min(a,1)^2\leq a$ and $(a-1)_+\leq(r-1)_+$,
\[
 \|a\|_2\leq\sqrt s+\epsilon_2.
\]
The first bracket of~\eqref{p1p:eq:decomposition} has half-variation at
most
\[
 \tfrac12\int_E(g+g\circ K)
 \leq\tfrac12\|g\|_2(2\sqrt s+\epsilon_2).
\]
The second has half-variation at most
$\frac12\|g\|_2\epsilon_2$. Lemma~\ref{p1p:lem:grid} gives
$\|g\|_2\leq H_z$ and total grid-approximation cost $V/(2N)$.
Integrating proves the result. In particular the argument controls
the image of the exceptional set even when $K$ merges inputs.
\end{proof}

This criterion has no ambient-dimensional exponent, but requires
stronger reference calibration. For a differentiable bijection of the
active square, the bound $|\det DK-1|\leq a<1$ implies
$|r-1|\leq a/(1-a)$ and hence $\epsilon_2\leq a/(1-a)$.
Ordinary reference TV does not supply this bound: compressing reference
mass $\epsilon$ into volume $\epsilon^2$ yields order-one $L^2$
distortion despite TV of order $\epsilon$.
Nor may the product in~\eqref{p1p:eq:sectiondist} be replaced by a product
of unweighted error averages. It is the actual section coefficient
$H_z$, which may concentrate, that weights the calibration error.

Finally, the necessity of a regularity input is already visible on the
unit circle. For $0<c<1$, take
$f_M(x)=1+c\cos(2\pi Mx)$ and translate by $1/(2M)$. The reference error
is zero, the displacement tends to zero, and the cap $1+c$ is fixed;
nevertheless the actual TV error is $2c/\pi$. Its variation is $4cM$.
Small reference error and small displacement therefore cannot replace
the law-sensitive regularity estimates in this section.


\section{One unknown input and a separate held receiver}
\label{p1i:sec:instrument}

The prepared writer can be used once to measure an unknown internal qubit.
The receiver need not have its wave population, and may remain correlated with
every old archive. What is needed is a quantitative restriction on that
complete law. The result below allows either a density cap or the physical-score
class of the preceding preparation theorem. It controls both the receiver's
faithfulness to the writer's own earlier declaration and the labelled
postmeasurement state of the input and a finite internal reference.

For equal-mass finite positive measures we use
$\operatorname{TV}(\mu,\nu)=\frac12\|\mu-\nu\|_1$; for equal-trace positive
operators we use $D(A,B)=\frac12\|A-B\|_1$. These conventions apply to the
subprobabilities used below. Let $\mathcal R$ be any finite-dimensional
internal Hilbert space, with no additional positional coordinate, and write
the unknown normalized input as
\begin{equation}\label{p1i:eq:input}
 \chi=|0\rangle\chi_0+|1\rangle\chi_1,
 \qquad p=\|\chi_0\|^2,\quad 1-p=\|\chi_1\|^2.
\end{equation}
No orthogonality of $\chi_0$ and $\chi_1$ in $\mathcal R$ is required.
The two displayed qubit sectors are orthogonal. A mixed input is covered by
a finite internal purification followed by a partial trace.

The entrance quantum wave is
\begin{equation}\label{p1i:eq:ready}
 \sqrt{g_\sigma(q)}\sqrt{h_a(z)}\,
 (|0\rangle\chi_0+|1\rangle\chi_1)\otimes\Xi(N),
\end{equation}
where $g_\sigma$ and $h_a$ are centered Gaussian densities with standard
deviations $\sigma$ and $a$. Here $q$ is the writer position, $z$ is a
\emph{different} receiver position, and $N$ includes every old archive and
genuinely retained external context. In \eqref{p1i:eq:ready}, $\Xi$ is
understood to depend on the positional part of $N$ and includes old
internal memories; a fixed external conditioning label need not have a
wave amplitude. The spectator holding dynamics must give zero current
in those old positional coordinates throughout this protocol, and
$\|\Xi\|>0$ on a set of full actual probability. These conditions hold
for the Gaussian bank with the holding projectors constructed above.
The actual entrance positional law is independent of
the choice of $\chi$, but is not asserted to equal the density of
\eqref{p1i:eq:ready}. An internal SWAP loads the unknown input only after the
known-input warmup. An input with an unknown positional reference already
entangled with these coordinates is outside \eqref{p1i:eq:ready}.

Let $F$ and $G$ be the CDFs of $g_\sigma$ and $h_a$, and put
$u=F(q_0)$ and $v=G(z_0)$. Quantiles are analysis coordinates only. The
Hamiltonian does not evaluate them or use the unknown number $p$ as a
control setting.
In this section $D$ and $A$ denote physical lengths; their dimensionless
values are $D/\sigma$ and $A/a$, respectively. These ratios are the
separation parameters denoted by $D,A$ in the standardized-coordinate
preparation analysis.

\subsection{The complete current and the generic-input map}

\begin{proposition}[Exact prescribed one-use parent]\label{p1i:prop:parent}
Use the canonical Schr\"odinger current for the complete configuration.
There is a finite spin-controlled oscillator protocol, with $C^\infty$ center
paths and time-dependent potential, that splits
the writer to $-D,+D$, copies into receiver packets at $-A,+A$, returns the
writer to its ready wave, and stores the measured internal qubit. Its
controls are the same for every input \eqref{p1i:eq:input}. The held receiver
has exactly zero complete current during return and subsequent storage.
For every finite hold interval on which the holding projectors and
spectator conditions are maintained, the decoder
\begin{equation}\label{p1i:eq:decoder}
 Y(z)=\begin{cases}0,&z<0,\\1,&z\geq0\end{cases}
\end{equation}
is fixed and time independent.

The active endpoint flow has the following description. Define
\begin{align}
 \rho_p(q)&=p g_\sigma(q+D)+(1-p)g_\sigma(q-D),\notag\\
 F_p(q)&=pF(q+D)+(1-p)F(q-D),\notag\\
 w_p(q)&=\frac{p g_\sigma(q+D)}{\rho_p(q)},\notag\\
 M_p(z)&=p h_a(z+A)+(1-p)h_a(z-A),\notag\\
 \widetilde w_p(z)&=\frac{p h_a(z+A)}{M_p(z)}.\label{p1i:eq:mixtures}
\end{align}
If $q$ denotes the split endpoint and $z$ the copy endpoint, then
\begin{align}
 u&=F_p(q),\label{p1i:eq:flow1}\\
 v&=w_p(q)G(z+A)+(1-w_p(q))G(z-A),\label{p1i:eq:flow2}\\
 r&=\widetilde w_p(z)F(q+D)+(1-\widetilde w_p(z))F(q-D),
 \label{p1i:eq:flow3}\\
 b&=pG(z+A)+(1-p)G(z-A).\label{p1i:eq:flow4}
\end{align}
Here $r=F(q_{\rm final})$. The map $T_p:(u,v)\mapsto(r,b)$ is a
measure-preserving bijection of the open unit square. All formulas include
$p=0,1$ without a singular branch-weight division.
\end{proposition}

\begin{proof}
For an oscillator of mass $m$ and frequency $\omega$, let $\phi$ be its real
normalized ground state, so $\sigma^2=\hbar/(2m\omega)$. For a prescribed
$C^3$ center $d(t)$, set
\begin{align}
 V_d(x,t)&=\tfrac12m\omega^2(x-d(t))^2-m\ddot d(t)x,\notag\\
 \psi_d(x,t)&=\phi(x-d(t))
  \exp\!\left(\frac{i m\dot d(t)x}{\hbar}+i\theta_d(t)\right),
 \qquad
 \dot\theta_d=-\frac{\omega}{2}-\frac{m\dot d^2}{2\hbar}.
 \label{p1i:eq:driven}
\end{align}
Substitution into the Schr\"odinger equation proves the formula: the
$\phi'$ terms match, the oscillator acts on $\phi$ with energy
$\hbar\omega/2$, and the acceleration term matches the time derivative
of $m\dot d x$. Its current is $j=\dot d|\psi_d|^2$.
The same construction applies to the receiver with its own mass,
frequency and width $a$.

Choose the flat smooth step \eqref{eq:smooth-step} on each finite stage
interval. The displayed Schr\"odinger calculation requires only a $C^3$
center, so the polynomial schedule given earlier is also sufficient for
the endpoint formulas and all estimates. In qubit sector $s$ use
the block-diagonal sum of \eqref{p1i:eq:driven} for the corresponding writer
and receiver centers. First move only the writer from $0$ to $(-1)^{s+1}D$;
then hold it and move only the receiver from $0$ to $(-1)^{s+1}A$;
then hold the receiver and return the writer to zero.

The two internal sectors eliminate every interference cross term in the
complete positional density and current, even when the reference vectors
overlap. Each moving-coordinate velocity is a convex combination of the
two prescribed center velocities. Active densities are strictly positive
at every finite coordinate, including $p=0,1$, and the velocities are
smooth locally and bounded on each finite time interval. Thus the active
ODE has a unique global solution for every finite entrance coordinate.
The old coordinates remain under their admitted held parent. During copy
the writer has zero current; during return the receiver has zero current.

For a one-dimensional conservative flow, its density CDF is constant
along a trajectory: differentiating the CDF in time gives $-j$, and the
advective term is $\rho\dot q=j$. Applying this identity first to the
writer mixture, then to the receiver at the frozen writer coordinate,
and finally to the returning writer at the frozen receiver coordinate
gives \eqref{p1i:eq:flow1}--\eqref{p1i:eq:flow3}. All the CDFs involved are
strictly increasing. The final joint wave is
\begin{equation}\label{p1i:eq:finalwave}
 \sqrt{g_\sigma(q_{\rm final})}
 \left(\sqrt{h_a(z+A)}e^{i\theta_0}|0\rangle\chi_0
 +\sqrt{h_a(z-A)}e^{i\theta_1}|1\rangle\chi_1\right)
 \otimes\Xi(N),
\end{equation}
with known spatially constant branch phases. Its receiver marginal is
$M_p$ and hence its receiver reference quantile is \eqref{p1i:eq:flow4}.

For a direct Jacobian check, write
\[
 L(q,z)=p g_\sigma(q+D)h_a(z+A)
       +(1-p)g_\sigma(q-D)h_a(z-A).
\]
It factors both as $\rho_p(q)$ times the conditional receiver density and
as $M_p(z)$ times the conditional split-writer density. Consequently the
Jacobian of $(u,v)$ with respect to $(q,z)$ is $L$, and that of $(r,b)$
with respect to $(q,z)$ is also $L$. Positivity gives determinant one for
$T_p$ and the successive inverse CDFs give its inverse. This proves
reference-area preservation independently of any actual-law assertion.

Finally let $U(t)$ be a spatially constant internal unitary that stores
the measured qubit in a fresh internal register. For an explicit SWAP,
take $K=\pi(I-\mathrm{SWAP})/2$ and
$U(t)=\exp[-i s(t)K]$, where $s$ increases smoothly from zero to one.
The transformed parent is
\begin{equation}\label{p1i:eq:swap}
 H_U(t)=U(t)H_{\rm hold}(t)U(t)^\dagger
                +i\hbar\dot U(t)U(t)^\dagger.
\end{equation}
The receiver holding projectors must be conjugated in this expression.
Then $U\Psi$ is the exact evolved wave. Spatial constancy makes its
complete density and canonical current identical to those of $\Psi$.
In particular, $j_z=0$ is preserved through the SWAP. The statement for a
later finite hold requires subsequent controls to preserve the stored
sector and receiver hold; it does not cover arbitrary future couplings
to that receiver. No unknown state is cloned or erased.
\end{proof}

\subsection{Own-outcome copying uniformly in the input}

Let $S$ be the sign of the writer at the end of its split, using the same
left/right convention as \eqref{p1i:eq:decoder}. This is an earlier actual
event. Define the whole-hold failure event
\[
 \mathcal F=\{\text{there exists }t\in[t_{\rm copy},T_H]
                         \text{ with }Y(z(t))\ne S\}.
\]
By Proposition~\ref{p1i:prop:parent}, $\mathcal F$ is exactly the endpoint
copy-mismatch event. Put $d_*=D/\sigma$, $a_*=A/a$ and let $\Phi$ be the
standard Gaussian CDF.

\begin{lemma}[Copying with a guarded receiver]\label{p1i:lem:copy}
Suppose the entrance law is $du\,\nu(dN,dv)$, of total mass $m$, and $\nu$
is supported on $\eta\leq v\leq1-\eta$. Choose a writer guard $c_w>0$,
$0<\gamma\leq\eta/2$, and $c_z<a_*$ with
$\Phi(-c_z)<\eta/2$. Then, uniformly in $p\in[0,1]$,
\begin{align}
 (du\,\nu)(\mathcal F)&\leq m B,\notag\\
 B&=\Phi(-(d_*-c_w))
                  +\frac{2\Phi(-(d_*+c_w))}{\gamma},
 \label{p1i:eq:Bcopy}\\
 \operatorname{TV}\bigl(\mathcal L(Y),m(p,1-p)\bigr)
                 &\leq m\{B+\Phi(-d_*)\}.\label{p1i:eq:outcome}
\end{align}
Here the label law has mass $m$. If instead $m=1$ and the actual receiver
marginal obeys $\nu(dN,dv)\leq C\kappa(dN)dv$ for a probability $\kappa$,
the same conclusions hold on replacing $B$ by $B+2C\eta$.
\end{lemma}

\begin{proof}
The split writer has density $\rho_p$. Its band $|q|<c_w\sigma$ has mass
at most $\Phi(-(d_*-c_w))$, since the two translated Gaussians give the
same band integral. On $q<-c_w\sigma$, the integral of the wrong posterior
is
\[
 \int_{q<-c_w\sigma}(1-w_p(q))\rho_p(q)\,dq
       =(1-p)\Phi(-(d_*+c_w)).
\]
On the right the corresponding integral is
$p\Phi(-(d_*+c_w))$. Markov's inequality bounds the two bad-posterior
sets by at most the second term in \eqref{p1i:eq:Bcopy}; the factor two is
a convenient conservative bound. There is no division by $p$ or $1-p$.

On the remaining left region, $w_p\geq1-\gamma$. The copy equation gives
\[
 G(z+A)\leq\frac{v}{w_p}
 \leq\frac{1-\eta}{1-\gamma}\leq1-\eta/2.
\]
Since $\Phi(-c_z)<\eta/2$, this implies
$z<(c_z-a_*)a<0$. The reflected argument gives $z>0$ on the right good
region. Independence of $u$ from the full $\nu$ means that the discarded
writer mass is multiplied by $m$, regardless of correlations between
$N$ and $v$. This proves the failure bound. Moreover
\[
 (du\,\nu)(S=0)=m\{p+(1-2p)\Phi(-d_*)\}.
\]
Coupling $Y$ with this earlier $S$ proves \eqref{p1i:eq:outcome}.
Under the cap, the omitted receiver-end intervals have mass at most
$2C\eta$, which proves the last statement.
\end{proof}

\subsection{The actual conditional-state metric}

At a nonnode, divide the complete wave by its positional norm to obtain
its internal conditional vector. For the parent above, tracing the old
factor and ready writer leaves the normalized logical-memory/reference
vector
\begin{equation}\label{p1i:eq:conditional}
 \xi_\chi(z)=
 \frac{\sqrt{h_a(z+A)}e^{i\theta_0}|0\rangle\chi_0
       +\sqrt{h_a(z-A)}e^{i\theta_1}|1\rangle\chi_1}
      {\sqrt{M_p(z)}}.
\end{equation}
For an actual endpoint law $\mu_{\chi,T}$ define the labelled state
\begin{equation}\label{p1i:eq:cq}
 \Omega_\mu(\chi)=\int
 |Y(z)\rangle\langle Y(z)|\otimes
 |\xi_\chi(z)\rangle\langle\xi_\chi(z)|\,d\mu_{\chi,T}.
\end{equation}
The target projective instrument is the trace-one operator
\begin{equation}\label{p1i:eq:ideal}
 \Lambda(\chi)=\sum_{s=0}^1|s\rangle\langle s|_Y
               \otimes|s\chi_s\rangle\langle s\chi_s|,
 \qquad |s\chi_s\rangle=|s\rangle\otimes\chi_s.
\end{equation}
The phases in \eqref{p1i:eq:conditional} do not change an individual ideal
branch projector. The classical register in \eqref{p1i:eq:cq} is an
analysis of the physical receiver coordinate; its introduction alone
would not construct a physical receiver.

For each fixed input the same-parent flow is a measurable map, so
entrance TV contracts under its pushforward. Integrating any trace-one
positive operator field contracts TV into trace distance:
\begin{equation}\label{p1i:eq:kernelcontract}
 D\left(\int K\,d\mu,\int K\,d\nu\right)
 \leq\tfrac12\int\|K\|_1\,d|\mu-\nu|
 =\operatorname{TV}(\mu,\nu).
\end{equation}
This elementary fact applies to \eqref{p1i:eq:cq}; it requires no
assumption that the non-Born input-to-output assignment is linear.

\begin{theorem}[One-use instrument from partial readiness]
\label{p1i:thm:instrument}
In the exact parent of Proposition~\ref{p1i:prop:parent}, the following
two statements hold uniformly for all inputs \eqref{p1i:eq:input} and all
finite internal reference dimensions.

\emph{Capped receiver.} Suppose
\[
 \operatorname{TV}(\mu,du\,\nu)\leq\delta,
 \qquad \nu(dN,dv)\leq C\kappa(dN)dv,
\]
where $\nu$ is the actual nuisance/receiver marginal and $\kappa$ is a
probability measure consistent with the retained spectators. With $B$ as
in \eqref{p1i:eq:Bcopy},
\begin{align}
 \Pr_\mu(\mathcal F)&\leq\delta+B+2C\eta,\notag\\
 D(\Omega_\mu(\chi),\Lambda(\chi))
 &\leq\delta+B+2C\eta+\Phi(-d_*)
                         +\sqrt{C\Phi(-a_*)}.\label{p1i:eq:capinstrument}
\end{align}

\emph{Common initial box, without a density cap.} Suppose the entrance
law decomposes as $\mu=\mu_B+\mu_E$, with masses $m$ and $\tau=1-m$.
The measure $\mu_B$ is the pushforward of a single original preparation
box, and $\nu_B$ is its own nuisance/receiver marginal. Assume
\begin{equation}\label{p1i:eq:boxpremises}
 \operatorname{TV}(\mu_B,du\,\nu_B)\leq\Delta_B,
 \qquad\nu_B(dN,dv)=f(N,v)\,\lambda(dN)dv,
 \qquad\int\operatorname{Var}_v f(N,\cdot)\,\lambda(dN)\leq V_z.
\end{equation}
Here $\lambda$ may include an arbitrary actual external-context law,
$f\geq0$, and $\nu_B$ is supported on $\eta\leq v\leq1-\eta$.
The one-use bounds are
\begin{align}
 \Pr_\mu(\mathcal F)&\leq\tau+\Delta_B+mB,\label{p1i:eq:boxhistory}\\
 D(\Omega_\mu(\chi),\Lambda(\chi))
 &\leq\tau+\Delta_B+m\{B+\Phi(-d_*)\}
       +(m+V_z/2)\sqrt{\Phi(-a_*)}.\label{p1i:eq:boxinstrument}
\end{align}
No claim of fresh reuse after this unknown-input measurement is included.
\end{theorem}

\begin{proof}
For a nonzero ideal branch weight let
$\eta_s=|s\chi_s\rangle/\|\chi_s\|$. For a zero-weight branch choose any
unit vector in its qubit sector. On the physical region $Y(z)=s$, the
pure-state trace distance from \eqref{p1i:eq:conditional} to $\eta_s$ is
$\sqrt{W(z)}$, where
\[
 W(z)=\begin{cases}
 (1-p)h_a(z-A)/M_p(z),&z<0,\\
 p h_a(z+A)/M_p(z),&z\geq0.
 \end{cases}
\]
This follows from the squared overlap with the orthogonal correct
qubit sector. It also holds for a zero ideal branch weight: the
conditional vector in that region is orthogonal to the chosen comparison
sector. Direct integration gives the exact identity
\begin{equation}\label{p1i:eq:wrongidentity}
 \int W(z)M_p(z)\,dz
 =(1-p)\int_{z<0}h_a(z-A)\,dz
       +p\int_{z\geq0}h_a(z+A)\,dz
 =\Phi(-a_*).
\end{equation}
The returned writer integrates to one. Area preservation of $T_p$ thus
gives
\begin{equation}\label{p1i:eq:pullwrong}
 \int_{(0,1)^2} W(z(T_p(u,v)))\,du\,dv=\Phi(-a_*).
\end{equation}

In the capped case, the comparison density is dominated by
$C\,du\,dv\,\kappa(dN)$. Equation~\eqref{p1i:eq:pullwrong} implies
$\mathbb E W\leq C\Phi(-a_*)$, and Cauchy gives a conditional-state
replacement cost at most $\sqrt{C\Phi(-a_*)}$. After replacement, only the
classical branch weights differ from \eqref{p1i:eq:ideal}; their trace
distance is the outcome TV from Lemma~\ref{p1i:lem:copy}. Finally
\eqref{p1i:eq:kernelcontract} adds $\delta$ once for replacement of the
actual entrance law by $du\,\nu$. The copy-history bound follows from the
same entrance-law replacement applied to the path event $\mathcal F$.

For the boxed assertion fix a spectator value $N$ and set
$m_N=\int f(N,v)dv$ and $V_N=\operatorname{Var}_v f(N,\cdot)$.
The one-dimensional BV inequality $f(N,v)\leq m_N+V_N$ almost everywhere
gives
\[
 \|f(N,\cdot)\|_2^2\leq m_N(m_N+V_N),
 \qquad \|f(N,\cdot)\|_2\leq m_N+V_N/2.
\]
This is equally the active-square $L^2(du\,dv)$ bound, since the comparison
density is independent of $u$. Cauchy and \eqref{p1i:eq:pullwrong}, followed
by integration in $N$, bound the entire subprobability state-replacement
cost by
\begin{equation}\label{p1i:eq:BVstatefee}
 \int f(N,v)\sqrt{W(z(T_p(u,v)))}\,du\,dv\,\lambda(dN)
 \leq(m+V_z/2)\sqrt{\Phi(-a_*)}.
\end{equation}
There is no need to normalize rare spectator conditionals or to assign
them a common cap. The outcome cost for this comparison law is
$m(B+\Phi(-d_*))$. The boxed actual entrance law differs by at most
$\Delta_B$, so \eqref{p1i:eq:kernelcontract} bounds its instrument
distance from $m\Lambda(\chi)$ by the last three terms in
\eqref{p1i:eq:boxinstrument}. The outside-box output and
$\tau\Lambda(\chi)$ are positive operators of the same trace $\tau$;
their trace distance is at most $\tau$. Adding them charges the original
tail once. The same decomposition gives \eqref{p1i:eq:boxhistory}, with
outside-box paths charged at most their mass. This proves both claims.
\end{proof}

\begin{remark}[Why a receiver premise is needed]\label{p1i:rem:receiver}
Writer readiness alone does not force an arbitrary actual receiver to
copy the writer. Fix $0<p<1$ and a finite left split-writer point $q$.
Since $0<w_p(q)<1$, the right-hand side of \eqref{p1i:eq:flow2} evaluated
at $z=0$ is strictly less than one. Choosing $v$ above that value gives
$z>0$ although $S=0$. By continuity the failure persists on a neighborhood
of $(q,v)$, on which a smooth actual density may be concentrated. This
explains the load-bearing cap or source-variation/support premise.
It also explains why the future receiver is included in the complete
preparation law rather than introduced later as an unpriced fresh stock.
\end{remark}

\subsection{The coherent 102-coordinate row}

\begin{corollary}[One-use consequence of the physical-score row]
\label{p1i:cor:PM100}
Take the preparation parent and complete actual-law class of
Theorem~\ref{p1p:thm:prep} with $n=100$ known-input warmup cycles.
The $102$ scalar positional coordinates are one writer, $100$ warmup
receivers, and one untouched receiver for the unknown measurement.
Assume the sum of the full conditional physical relative scores is at
most $100$, and the sum of actual standardized twentieth moments is at
most $10^{15}$. All these conditions include the untouched receiver and
are averaged against the actual retained external-context law. Use the
complete initial box $|\xi_i|\leq10$ and separations $50$ packet widths.
Load the unknown input after warmup and use the same separation for its
one-use measurement. Then
\begin{equation}\label{p1i:eq:headline}
 \sup_{\chi,\,\dim\mathcal R<\infty}
 D(\Omega_\mu(\chi),\Lambda(\chi))<1.0052\times10^{-5}.
\end{equation}
The same ceiling bounds the unknown measurement's own-outcome copy and
entire admitted hold failure. It does not assert that the unknown-use
output is fresh for another unknown measurement.
\end{corollary}

\begin{proof}
The original complete box has outside mass
$\tau\leq10^{15}/10^{20}=10^{-5}$. The preparation estimates give
\begin{align}
 V_w&\leq2.6331\times10^{23},&
 E_{\rm map}&\leq2.00006\times10^{-16},\notag\\
 \Delta_B&\leq\frac{V_w}{4\,2^{100}}+2E_{\rm map}
                  <5.192874163845\times10^{-8}.\label{p1i:eq:boxrow}
\end{align}
The quantity $\Delta_B$ compares the actual \emph{boxed} warmup endpoint
to a uniform writer times its own boxed nuisance marginal. It does not
already include the outside-box tail.

The unused receiver coordinate $v$ does not participate in warmup.
Each warmup map is independent of $v$ and preserves the reference volume
of the other coordinates. Therefore its pushforward preserves the
integrated $v$-directional BV norm of the initial boxed density; this can
be seen by commuting the distributional $v$ derivative with the
pushforward and using volume preservation. Marginalizing the final
writer contracts that norm. The physical score-to-BV estimate therefore
gives
\begin{equation}\label{p1i:eq:sourcevariation}
 V_z\leq\sqrt{2\pi}e^{50}
                \{\sqrt{J_z}+\sqrt{J_z+2}\}
        <2.6331\times10^{23},
\end{equation}
where $J_z\leq100$. The source remains supported on
$v\in[\Phi(-10),1-\Phi(-10)]$ under every warmup map and its
marginalization. Thus it already satisfies the receiver guard for
$\eta=10^{-110}$; no new source cut or second moment-tail charge is used.

Apply Theorem~\ref{p1i:thm:instrument} with
$d_*=a_*=50$, $c_w=3$, $c_z=24$ and $\gamma=10^{-130}$.
The inequalities $\gamma\leq\eta/2$ and $\Phi(-24)<\eta/2$ hold, and
the left good receiver finishes below $-26a$. Gaussian tail bounds give
\begin{align}
 B&=\Phi(-47)+2\Phi(-53)/10^{-130}<10^{-480}+2\times10^{-470},
 \notag\\
 \Phi(-50)&<10^{-544},\notag\\
 (1+V_z/2)\sqrt{\Phi(-50)}&<10^{-240}.
 \label{p1i:eq:smallfees}
\end{align}
For reproducibility, the elementary inequality
$\Phi(-x)<e^{-x^2/2}/(x\sqrt{2\pi})$, obtained by bounding
$1\leq t/x$ in the Gaussian tail integral, proves these comparisons.
The receiver guard follows also from
$\Phi(-10)\geq0.01\,\phi(10.01)>10^{-110}$, by integrating just over
$[10,10.01]$. Together with \eqref{p1i:eq:boxrow}, substitution in
\eqref{p1i:eq:boxinstrument} and \eqref{p1i:eq:boxhistory} yields the
strict ceiling \eqref{p1i:eq:headline}.
\end{proof}

The count $102$ concerns scalar configurational coordinates of this
effective oscillator parent. It is neither a count of microscopic
particles nor the dimension of the internal quantum space. The finite
internal reference and stored qubits introduce no new positional law
in this model. The common-box proof keeps the same original
subprobability throughout preparation and measurement. Its single tail
allowance accounts for all coordinates of the original bank, including
the future receiver; cutting that receiver again would unnecessarily
charge part of the same exceptional population twice.

\section{Instrument stability, reference systems, and postselection}
\label{p1i:sec:stability}

The exact-parent theorem also identifies what a perturbative
implementation must supply. Two complementary transfer statements are
useful: a full-ready operator-isometry criterion and a cap-free
conditional-state criterion. Neither converts an endpoint wave norm
alone into a displacement bound or a whole-history record theorem.

\subsection{An abstract full-ready criterion}

\begin{proposition}[Finite-reference instrument transfer]
\label{p1i:prop:abstract}
Let an input-independent ready positional law $\mu_0$ obey
$\operatorname{TV}(\mu_0,\rho_0)\leq\delta$. Suppose the exact physical
ready wave is an isometric product embedding of an internal input, with
the same positional reference $\rho_0$ for every input. For every input
assume a specified conservative complete flow and a pointwise physical
wave representative on the actual-law support. Let $A$ be its endpoint
wave isometry and let
\[
 A_0\chi=\sum_s\phi_s\otimes V_s\chi,
 \qquad \|\phi_s\|=1,
 \qquad V_s^\dagger V_t=0\ (s\ne t),
 \qquad\sum_s V_s^\dagger V_s=I.
\]
The $V_s$ have orthogonal internal ranges; any retained branch-dependent
purifier belongs to $\phi_s$ and to its full positional integral.
Suppose $\|A-A_0\|_{\rm op}\leq\varepsilon$ and the disjoint record
regions $R_s$ partition configuration space with
$\int_{R_s^c}\|\phi_s(q)\|^2dq\leq t$, $0\leq t\leq1$.
Then the actual labelled conditional-state instrument satisfies
\begin{equation}\label{p1i:eq:abstract}
 D\!\left(\Omega_\mu(\chi),
       \sum_s|s\rangle\langle s|\otimes
            V_s|\chi\rangle\langle\chi|V_s^\dagger\right)
 \leq\min\{1,\delta+\varepsilon+\sqrt{2t-t^2}\}.
\end{equation}
The estimate extends to every finite internal reference with no
dimension multiplier. If the flow is only admitted on a reference-full
set, any actual mass outside its domain must instead be explicitly
charged as unresolved, rather than treated as a predicted trajectory.
\end{proposition}

\begin{proof}
For each fixed input, equivariance and TV contraction transport
$\delta$ to the endpoint. Equation~\eqref{p1i:eq:kernelcontract} therefore
reduces the proof to its wave-populated endpoint. Under that law the
normalization denominator in the conditional projector cancels with
the positional density, so the actual cq integral equals the usual
labelled wave extraction.

Define the isometry $B\psi=\sum_s|s\rangle 1_{R_s}\psi$ and the
comparison vector $C_0\chi=\sum_s|s\rangle\phi_s V_s\chi$. Put
$w_s=\|V_s\chi\|^2$ and
$L=\sum_s w_s\int_{R_s^c}\|\phi_s\|^2\leq t$.
Orthogonality of the branch ranges gives exactly
\[
 \langle BA_0\chi,C_0\chi\rangle=1-L,
 \qquad \|BA_0\chi-C_0\chi\|^2=2L.
\]
Hence their pure-state trace distance is
$\sqrt{1-(1-L)^2}\leq\sqrt{2t-t^2}$.
Apply the same extraction channel to both vectors: trace all positional
and discarded purifier degrees and \emph{dephase the classical label}.
For $BA\chi$ the disjoint regions already make the extracted label
diagonal. For $C_0\chi$, dephasing is essential when the $\phi_s$ overlap;
after it the result is precisely the ideal instrument in
\eqref{p1i:eq:abstract}. Trace distance contracts under this channel.
The remaining wave error is at most $\varepsilon$, since the distance
of normalized pure states is bounded by their vector-norm difference
and $B$ is an isometry.

For a finite reference $\mathcal R$, decompose an arbitrary vector in an
orthonormal reference basis. Summing squared output norms proves
$\|(A-A_0)\otimes I_{\mathcal R}\|\leq\|A-A_0\|$; product vectors give
the reverse inequality. All other steps retain the reference and are
unchanged. Convexity and purification give the mixed-input statement.
\end{proof}

An operator bound in this proposition is stronger than the maximum of
the errors on a chosen input basis. If the basis-column errors are
$e_j$, the justified bound is $\varepsilon\leq(\sum_j e_j^2)^{1/2}$,
by Cauchy--Schwarz, unless a sharper operator estimate is supplied.
Likewise a position-dependent output-frame rotation must be included in
$A_0$ or charged in $\varepsilon$. The theorem compares complete
conditional states, so an agreement of position densities alone is
insufficient.

\subsection{Cap-free transfer of conditional quantum states}

\begin{proposition}[State stability under physical-score concentration]
\label{p1i:prop:capfreestate}
Let $d\geq2$ and let $X$ be retained external context with its actual
law. Conditional on $X$, suppose the initial boxed actual density
$f_{B,X}$ relative to the normalized physical entrance reference obeys
\[
 \|f_{B,X}\|_{d/(d-1)}\leq B_X,
 \qquad\|B_X\|_{L^2(\mathrm{actual}\ X)}\leq B_2.
\]
For example, the physical-score BV bound gives
$B_X=m_X+V_{{\rm total},X}/d$ and
\[
 B_2\leq1+\frac{D_R}{\sqrt d}
                  \{\sqrt J+\sqrt{J+2d}\}.
\]
Let $\Psi_X$ and $\Psi_{0,X}$ be normalized complete endpoint waves on
the same physical configuration space, retaining the input, internal
reference and purifiers. The physical flow transports the entrance
reference to $\rho_X=\|\Psi_X\|^2$ and the boxed actual law to its actual
endpoint. Set
\[
 e_X=\|\Psi_X-\Psi_{0,X}\|_2,
 \qquad e_{\rm rms}^2=\mathbb E_{\mathrm{actual}\ X} e_X^2.
\]
Replacing the physical conditional-state field by the ideal one in the
boxed actual cq output costs at most
\begin{equation}\label{p1i:eq:capfreestate}
 s_B=B_2 e_{\rm rms}^{2/d}.
\end{equation}
If a complete boxed endpoint-law comparison separately costs $\ell_B$
and the ideal boxed instrument differs from $m\Lambda$ by at most
$\eta_{0,B}$, then the full instrument satisfies
\begin{equation}\label{p1i:eq:physicalcomposition}
 D(\Omega_{\rm physical},\Lambda)
                   \leq\tau+\eta_{0,B}+\ell_B+s_B.
\end{equation}
The wave and law hypotheses must hold uniformly over the allowed
internal inputs and references for this to be a uniform instrument
statement.
\end{proposition}

\begin{proof}
At a physical nonnode let $a=\Psi_X(q)$ and, at a nonzero ideal vector
$b=\Psi_{0,X}(q)$, let $P_b$ be its rank-one orthogonal projector. If
$D_q$ is the trace distance between their normalized pure states, then
\[
 \rho_X(q)D_q^2=\|(I-P_b)a\|^2
                         \leq\|a-b\|^2.
\]
At an ideal node choose any unit comparison vector; $D_q\leq1$ gives
the same inequality. Thus $\int D_q^2\rho_X\,dq\leq e_X^2$.

The relative density of the boxed actual endpoint with respect to
$\rho_X$ is the conditional expectation of $f_{B,X}$ on the physical
endpoint map. Jensen therefore contracts its $L^{d/(d-1)}$ norm; if the
flow is invertible the norm is preserved. H\"older's inequality and
$0\leq D_q\leq1$ give
\[
 \int D_q\,d\mu_{B,X,T}
 \leq B_X\left(\int D_q^d\rho_X\,dq\right)^{1/d}
 \leq B_X e_X^{2/d}.
\]
Cauchy in the actual $X$ law, followed by concavity of $x^{2/d}$, yields
\[
 \mathbb E B_Xe_X^{2/d}
 \leq B_2(\mathbb E e_X^{4/d})^{1/2}
 \leq B_2(\mathbb E e_X^2)^{1/d}.
\]
Copying a common endpoint label and tracing unwanted internal degrees
can only decrease this fee. Applying
\eqref{p1i:eq:kernelcontract} to the separately certified boxed law
comparison adds $\ell_B$; comparing the ideal boxed output to
$m\Lambda$ adds $\eta_{0,B}$. The outside-box output and $\tau\Lambda$
have the same trace and cost at most $\tau$. This proves
\eqref{p1i:eq:physicalcomposition}.
\end{proof}

The norm $e_{\rm rms}$ is weighted by the \emph{actual} retained context.
A small wave norm averaged against an unrelated quantum distribution of
$X$ does not supply this hypothesis. A context of arbitrarily small
quantum weight can carry all the actual mass and an order-one state
error. Also the physical versus ideal endpoint-law comparison in
\eqref{p1i:eq:physicalcomposition} must include every drifting position;
one cannot retain it merely as an unpriced ``spectator.''

For the $d=102$ row, $B_2<4\times10^{22}$ and
$\eta_{0,B}<5.2\times10^{-8}$. Thus separately established bounds
$\ell_B\leq2.8\times10^{-8}$ and $e_{\rm rms}\leq10^{-3200}$ would give
$s_B<4\times10^{-40}$ and a full instrument error below
$1.0081\times10^{-5}$. These are sufficient finite interface
requirements. They are not measured calibration tolerances or a claim
that a material device achieves them. In particular they do not follow
from the exact prescribed oscillator solution.

A perturbed whole-hold claim needs a path estimate of its own. If the
complete physical current proves that the possibly failing entrance
set has conditional reference volume $b_X$, the same H\"older--Cauchy
argument, now applied to its indicator, gives
\begin{equation}\label{p1i:eq:pathfee}
 \Pr_{\rm actual}(\text{path failure})
 \leq\tau+B_2(\mathbb E_{\mathrm{actual}\ X}b_X)^{1/d}.
\end{equation}
This is additional information, not a consequence of
\eqref{p1i:eq:physicalcomposition}.

\subsection{What the cq estimate says after selecting an outcome}

\begin{proposition}[Postselection bound]\label{p1i:prop:postselection}
Let two normalized cq states have blocks $a_s\sigma_s$ and $p_s\tau_s$
and satisfy $D(\Omega,\Lambda)\leq\epsilon$. Then
\begin{equation}\label{p1i:eq:postselection}
 \sum_s a_sD(\sigma_s,\tau_s)\leq2\epsilon,
 \qquad
 D(\sigma_s,\tau_s)\leq\min\{1,\epsilon/a_s\}\quad(a_s>0).
\end{equation}
There is no uniform normalized-state conclusion for arbitrarily rare
actual outcomes.
\end{proposition}

\begin{proof}
Write $\Delta_s=a_s\sigma_s-p_s\tau_s$ and assign arbitrary comparison
states when $p_s=0$. The triangle inequality gives
\[
 a_s\|\sigma_s-\tau_s\|_1
             \leq\|\Delta_s\|_1+|a_s-p_s|.
\]
Summing and using contraction to the classical label proves the first
inequality. For an individual outcome, the complementary blocks have
total trace $-(a_s-p_s)$, so their summed trace norms are at least
$|a_s-p_s|$. Hence
$\|\Delta_s\|_1+|a_s-p_s|\leq2\epsilon$, proving the second inequality.
For sharpness, take one actual outcome with probability $0.02$ and
conditional state $|0\rangle$, its ideal probability $0.01$ and state
$|1\rangle$, and identical conditional states in the complementary
outcome. The unconditional cq distance is $0.02$, but the selected
conditional states have distance one.
\end{proof}

All these bounds are uniform entangled-input comparisons within the
stated internal stock. They do not identify an exactly linear completely
positive actual map for a non-Born law, and hence are not assertions of
a diamond norm for such a map. The exact-parent theorem establishes one
unknown use after preparation. A repeated-use theorem would additionally
have to retain the relevant contexts, establish fresh stocks at each
stage, and specify which retained information may control subsequent
ideal operations.

\subsection{Physical scope of the one-use result}

The positive result is an explicit finite canonical-current parent with
a complete-law preparation bound and a uniform internal-reference
instrument. Its physical premises are specific: Gaussian quantum
stocks, spin-dependent translated traps including the inertial term in
\eqref{p1i:eq:driven}, spatially constant internal rotations, correctly
conjugated holding projectors, and the stated full actual-law score and
moment bounds. The result does not derive those actual-law conditions
from ordinary wave energy or from the factorized quantum stock.

A microscopic implementation must separately establish the full
Hamiltonian/current reduction, coordinate conventions, finite source
and clock controls, their retained positional states, uniform-input
wave and map errors, and the required held-history bound. These are
coefficient and modeling obligations on a proposed implementation, not
unstated hypotheses absorbed by the small error in
\eqref{p1i:eq:headline}. Controlled breathing-Gaussian alternatives can
provide model-specific perturbation examples, but their frequency,
width, receiver and actual-law calibration assumptions must be proved
compatible before their constants can be combined with this
102-coordinate physical-score row.

\appendix
\section{A complementary finite regular-basin theorem}\label{app:regular}
This appendix concerns a different, random-return preparation mechanism. It is not used to manufacture independent archive populations for the deterministic Gaussian protocol. Its purpose is to distinguish uniform preparation of a regularity class from retained-tape independence, and to record the supporting finite-dimensional argument.

Let $\gamma$ be the standard Gaussian measure on $\mathbb R^d$. Let $P$ be a self-adjoint Markov contraction on $L^2(\gamma)$, preserving constants, positive semidefinite as an operator, and with no nonconstant fixed function. The spectral positivity excludes a period-two obstruction. Since $P$ preserves nonnegative functions and $P1=1$, self-adjointness gives $\int Ph\,d\gamma=\int h\,d\gamma$ and
\[
 \|Ph\|_1\leq\int P|h|\,d\gamma=\|h\|_1
 \qquad(h\in L^2(\gamma)).
\]
Thus $P$ extends uniquely by density to a positivity- and mass-preserving
contraction on $L^1(\gamma)$. All expressions $P^k f$ below use this
extension; the admitted density $f$ need not lie in $L^2$.
A lazy inverse-balanced mixture of actual reference-preserving returns
is one possible supplier; its physical implementation and command
independence are separate hypotheses. Denote normalized product Hermite
polynomials by $H_\alpha$, and set
\[
 E_m=\operatorname{span}\{H_\alpha:|\alpha|\le m\},\quad
 S_m=\sum_{|\alpha|\le m}H_\alpha^2,\quad C_m^2=\int S_m^2\,d\gamma.
\]
For $E_\ell^0=E_\ell\cap1^\perp$, define $\kappa_\ell(k)=\lVert P^k|_{E_\ell^0}\rVert_{L^2\to L^2}$, taking it to be zero if the domain is zero dimensional. The range need not be $E_\ell^0$.

\begin{theorem}[Finite preparation on a weak Fisher basin]\label{thm:hermite}
If $f\ge0$, $\int f\,d\gamma=1$, $\sqrt f\in H^1(\gamma)$, and $4\int|\nabla\sqrt f|^2\,d\gamma\le J$, then, for $\delta_m=\sqrt{J/[4(m+1)]}<1$,
\[
 \TV\bigl((P^k f)\gamma,\gamma\bigr)
 \le\delta_m+\frac12\sqrt{C_m^2-1}\,\kappa_{2m}(k).
\]
For every $J<\infty$ and $\epsilon>0$ there is a finite common $k$ for the entire stated class. No spectral gap or practical waiting-time bound is asserted.
\end{theorem}
\begin{proof}
Put $g=\sqrt f$, $h=\Pi_mg$, and $c=\lVert h\rVert_2^2$. Gaussian integration by parts gives
\[
 \langle\partial_i g,H_\beta\rangle
 =\sqrt{\beta_i+1}\,\langle g,H_{\beta+e_i}\rangle.
\]
The identity extends from smooth functions to $H^1(\gamma)$ by Sobolev
approximation. Bessel's inequality, summed over $i$, gives
$\sum_\alpha|\alpha||\langle g,H_\alpha\rangle|^2
 \leq\int|\nabla g|^2\,d\gamma\leq J/4$. Hence $1-c=\lVert g-h\rVert_2^2\le\delta_m^2<1$. The polynomial $a=h^2/c$ is a normalized positive density. Since $g\ge0$,
\[
 \int g|h|/\sqrt c\,d\gamma\ge\langle g,h\rangle/\sqrt c=\sqrt c.
\]
For normalized nonnegative $r,s$, Cauchy--Schwarz applied to $(r-s)(r+s)$ gives $\frac12\int|r^2-s^2|\le\sqrt{1-\langle r,s\rangle^2}$. Thus $\TV(f\gamma,a\gamma)\le\sqrt{1-c}\le\delta_m$, even if $h$ changes sign. Pointwise coefficient Cauchy--Schwarz gives $a\le S_m$, whence $\lVert a-1\rVert_2^2\le C_m^2-1$. The polynomial $a-1$ is in $E_{2m}^0$. Markov contraction and $L^1\le L^2$ prove the bound.

The spectral theorem gives $P^kg\to0$ for each $g\perp1$: the spectrum lies in $[0,1]$, and the spectral mass at $1$ is absent. Convergence is uniform on the unit sphere of any fixed finite-dimensional domain, so $\kappa_{2m}(k)\to0$. Choose $m$ first to make $\delta_m<\epsilon/2$, and then $k$ for the second term. This order avoids both a spectral-gap assumption and an unsupported invariance assumption on the Hermite subspace.
\end{proof}

The finite Gram matrix $G_{\ell,k}(\alpha,\beta)=\langle H_\alpha,P^{2k}H_\beta\rangle$ has largest eigenvalue $\kappa_\ell(k)^2$. Its trace is an upper bound. For $m=1$, $S_1(x)=1+|x|^2$ and the Gaussian moments $\mathbb E|x|^2=d$, $\mathbb E|x|^4=d(d+2)$ give $C_1^2=1+d^2+4d$. For $d=1$ and $m=2$, $S_2(x)=(3+x^4)/2$, so $C_2^2=(9+6\cdot3+105)/4=33$. A numerical waiting time for a concrete nonlinear return library would additionally require enclosed Gram entries; the theorem does not assert that this calculation has been completed.

\begin{proposition}[Retained tape and actual correlations]\label{prop:tape}
Suppose $W$ is an independent command word with law $\pi$, and every $F_w$ is an invertible measurable $\gamma$-preserving map. Then
\[
 \TV\bigl(\operatorname{Law}(F_WZ,W),\gamma\otimes\pi\bigr)
 =\TV\bigl(\operatorname{Law}(Z),\gamma\bigr).
\]
If only the last $r$ commands of a fresh independent sequence of $k$ commands are retained, the corresponding discrepancy equals that of the marginal configuration before those last $r$ commands.
\end{proposition}
\begin{proof}
The measurable bijection $(z,w)\mapsto(F_wz,w)$ sends $\gamma\otimes\pi$ to itself. Total variation is invariant under a common measurable bijection. For a retained suffix, its independence from the preceding configuration gives the same product input argument starting at time $k-r$.
\end{proof}
This is the total-variation version of the inherited retained-memory principle, not a claim that information conservation is new. Even a one-bit archive may retain the entire relevant map: if commands apply either the identity or an involution $R$, the parity of the number of $R$ commands determines the endpoint map and permits its inverse echo. Correct command marginals alone are also insufficient. With an input density $a(F_wz)$ relative to $\gamma(dz)\pi(dw)$, each command marginal is still $\pi$, while the endpoint density is exactly $a$. Taking $a=2\mathbf1_H$ for a reference half-space $H$ gives endpoint discrepancy $1/2$.

\begin{proposition}[Finite or countable command obstruction]\label{prop:atomic}
A finite or countable random mixture of invertible deterministic commands sends an initial point mass to a countably supported measure. Its total variation from a non-atomic Gaussian remains one, irrespective of the command probabilities.
\end{proposition}
\begin{proof}
The set of reachable points is countable and has output probability one but Gaussian probability zero.
\end{proof}
Accordingly a continuous-parameter smoothing construction must explicitly supply its analog command law and a rank condition; it is a different statistical resource. The present finite Gaussian protocol instead restricts the complete original actual law by weak scores. Neither construction obtains its required law from the mere smoothness of a quantum wave.

\section{Physical scope and relation to companion work}\label{app:scope}
The central results refer to one exact controlled harmonic parent, one full original conditional law and one finite internal-reference instrument. The companion manuscript on repeated position records~\cite{RodgersP2} uses a separate effective scalar-current construction with retained entropy and calibrated reset. Its repeated-record conclusion does not follow merely by reusing the writer prepared here. The singular-flow analysis~\cite{RodgersP3} and coherent-source current estimates~\cite{RodgersP4} address possible interfaces for more general physical models; their hypotheses are not needed for the explicit positive Gaussian flows of this paper. Applying them to a material implementation would require a compatibility proof for the complete Hamiltonian, current, original law and decoder. These new companion manuscripts are separate from the current website research editions, and no completed integration with that programme is asserted. In particular, assumptions from its pilot-medium and massive-configuration constitutions are not transferred between them here.

An arbitrary positional reference is especially different from the finite internal reference used above. If an untouched positional spectator has actual law $\eta$ and quantum reference $\eta_0$, the joint output discrepancy is at least $\TV(\eta,\eta_0)$ by marginalization. Preparing only the apparatus cannot reduce this lower bound. Likewise, an exact initial archive coordinate copied into $X$ must remain in the conditioning; it may make the full conditional law singular. The purpose of these examples is to specify what the theorem actually prepares, rather than to deny the mathematical value of conditional preparation.

\section*{Conclusion}
Reversible finite Gaussian dynamics can prepare a writer conditionally on its complete retained archive for a broad, explicitly non-Born class of original laws. The exact smooth map is essential: fine archive information invalidates an unrestricted writer-only premise, while full conditional physical scores and actual moments yield constructive quantitative bounds without a density cap. The same parent supports a separate physical receiver and an input-uniform finite-internal-reference instrument. These are positive mathematical results with identifiable statistical and control premises. Establishing those premises independently in a material apparatus, and extending the argument to repeated unknown inputs and arbitrary positional references, remain distinct scientific tasks.

\section*{Author statement and research support}
This research was conducted independently by Jeremy Rodgers without external funding. AI systems assisted with drafting, mathematical derivations, source and bibliography checks, and computational verification. The author is responsible for the manuscript; AI-assisted checking does not constitute independent scientific validation. Collaboration on independent mathematical verification and physical implementation is welcome.

\begin{thebibliography}{99}
\raggedright
\bibitem{Rosenblatt1952} M. Rosenblatt, Remarks on a Multivariate Transformation, \emph{Ann. Math. Statist.} \textbf{23} (1952), 470--472. \href{https://doi.org/10.1214/aoms/1177729394}{doi:10.1214/aoms/1177729394}.
\bibitem{Gross1975} L. Gross, Logarithmic Sobolev Inequalities, \emph{Amer. J. Math.} \textbf{97} (1975), 1061--1083. \href{https://doi.org/10.2307/2373688}{doi:10.2307/2373688}.
\bibitem{LasotaYorke1973} A. Lasota and J. A. Yorke, On the Existence of Invariant Measures for Piecewise Monotonic Transformations, \emph{Trans. Amer. Math. Soc.} \textbf{186} (1973), 481--488. \href{https://doi.org/10.1090/S0002-9947-1973-0335758-1}{doi:10.1090/S0002-9947-1973-0335758-1}.
\bibitem{Bennett1973} C. H. Bennett, Logical Reversibility of Computation, \emph{IBM J. Res. Dev.} \textbf{17} (1973), 525--532. \href{https://doi.org/10.1147/rd.176.0525}{doi:10.1147/rd.176.0525}.
\bibitem{ValentiniWestman2005} A. Valentini and H. Westman, Dynamical Origin of Quantum Probabilities, \emph{Proc. R. Soc. A} \textbf{461} (2005), 253--272. \href{https://doi.org/10.1098/rspa.2004.1394}{doi:10.1098/rspa.2004.1394}. \href{https://arxiv.org/abs/quant-ph/0403034}{arXiv:quant-ph/0403034}.
\bibitem{GoldsteinStruyve2007} S. Goldstein and W. Struyve, On the Uniqueness of Quantum Equilibrium in Bohmian Mechanics, \emph{J. Stat. Phys.} \textbf{128} (2007), 1197--1209. \href{https://doi.org/10.1007/s10955-007-9354-5}{doi:10.1007/s10955-007-9354-5}. \href{https://arxiv.org/abs/0704.3070}{arXiv:0704.3070}.
\bibitem{DurrGoldsteinZanghi1992} D. D\"urr, S. Goldstein and N. Zangh\`i, Quantum Equilibrium and the Origin of Absolute Uncertainty, \emph{J. Stat. Phys.} \textbf{67} (1992), 843--907. \href{https://doi.org/10.1007/BF01049004}{doi:10.1007/BF01049004}. \href{https://arxiv.org/abs/quant-ph/0308039}{arXiv:quant-ph/0308039}.
\bibitem{DurrGoldsteinZanghi2004} D. D\"urr, S. Goldstein and N. Zangh\`i, Quantum Equilibrium and the Role of Operators as Observables in Quantum Theory, \emph{J. Stat. Phys.} \textbf{116} (2004), 959--1055. \href{https://doi.org/10.1023/B:JOSS.0000037234.80916.d0}{doi:10.1023/B:JOSS.0000037234.80916.d0}. \href{https://arxiv.org/abs/quant-ph/0308038}{arXiv:quant-ph/0308038}.
\bibitem{DurrEtAl2005Density} D. D\"urr, S. Goldstein, R. Tumulka and N. Zangh\`i, On the Role of Density Matrices in Bohmian Mechanics, \emph{Found. Phys.} \textbf{35} (2005), 449--467. \href{https://doi.org/10.1007/s10701-004-1983-9}{doi:10.1007/s10701-004-1983-9}. \href{https://arxiv.org/abs/quant-ph/0311127}{arXiv:quant-ph/0311127}.
\bibitem{DaviesLewis1970} E. B. Davies and J. T. Lewis, An Operational Approach to Quantum Probability, \emph{Commun. Math. Phys.} \textbf{17} (1970), 239--260. \href{https://doi.org/10.1007/BF01647093}{doi:10.1007/BF01647093}.
\bibitem{Ozawa1984} M. Ozawa, Quantum Measuring Processes of Continuous Observables, \emph{J. Math. Phys.} \textbf{25} (1984), 79--87. \href{https://doi.org/10.1063/1.526000}{doi:10.1063/1.526000}.
\bibitem{FuchsVanDeGraaf1999} C. A. Fuchs and J. van de Graaf, Cryptographic Distinguishability Measures for Quantum-Mechanical States, \emph{IEEE Trans. Inform. Theory} \textbf{45} (1999), 1216--1227. \href{https://doi.org/10.1109/18.761271}{doi:10.1109/18.761271}. \href{https://arxiv.org/abs/quant-ph/9712042}{arXiv:quant-ph/9712042}.
\bibitem{RodgersReturn2026} J. Rodgers, \emph{Preparation-Return Holonomy and Equilibrium Uniqueness in Engineered Interacting Networks}, version 2, 4 October 2026, publication revision of the 7 September 2026 manuscript. \href{https://doi.org/10.5281/zenodo.23131064}{doi:10.5281/zenodo.23131064}.
\bibitem{RodgersControl2026} J. Rodgers, \emph{Control Consistency and Born Equilibrium: Uniqueness from Local Potentials and Fixed Interactions}, version 2, revised preprint candidate, 4 October 2026. \href{https://doi.org/10.5281/zenodo.23131058}{doi:10.5281/zenodo.23131058}.
\bibitem{RodgersP2} J. Rodgers, \emph{Effective Repeated Position Records with Retained Entropy and Calibrated Reset}, companion manuscript, 2026. \href{https://doi.org/10.5281/zenodo.23259558}{doi:10.5281/zenodo.23259558}.
\bibitem{RodgersP3} J. Rodgers, \emph{Reference-Weighted Deterministic Quantum Flows with Singular Interactions and Retained Sources}, companion manuscript, 2026. \href{https://doi.org/10.5281/zenodo.23259569}{doi:10.5281/zenodo.23259569}.
\bibitem{RodgersP4} J. Rodgers, \emph{Complete-Current Estimates for Coherent Quantum Sources and Continua}, companion manuscript, 2026. \href{https://doi.org/10.5281/zenodo.23259571}{doi:10.5281/zenodo.23259571}.
\end{thebibliography}

\end{document}
